% 第 X 讲，第 6 节至本讲末；原 PDF 第 87—106 页。
\SourcePage{87}
\SGAsection{6}{无限 Galois 主覆盖与扩大基本群}
（本节与下一节的结果，在本研讨班后续部分不再使用。）

设 $S$ 是一个\Term{prescheme}{预概形}。我们打算确定 $S$ 上结构群形如 $G_S$ 的\Term{principal-homogeneous-bundle}{主齐性纤维丛} $P$；这里 $G_S$ 是普通群 $G$（不必\Term{finite}{有限}）所定义的 $S$-\Term{constant-group}{常值群}。这也称为 $S$ 上群为 $G$ 的 \Term{galois-principal-bundle}{Galois 主纤维丛}。我们按\Term{fpqc-topology}{忠实平坦拟紧拓扑}的意义理解“\Term{principal-bundle}{主纤维丛}”（参见第 IV 讲，定义 5.1.5）；但须指出，对这样的 $P$，结构\Term{morphism}{态射} $P\to S$ 必然\Term{etale}{平展}且\Term{surjective}{满射}，因而是\Term{etale-topology}{平展拓扑}的\Term{covering-morphism}{覆盖态射}，所以 $P$ 对\Term{etale-topology}{平展拓扑}也是\Term{locally-trivial}{局部平凡}的（参见 IV，命题 5.1.6）。\footnote{编者注（47）：注意这里假定 $P$ 是\Term{prescheme}{预概形}。在先验更一般的情形，即一个（fpqc）\Term{sheaf}{层} $P$ 为 $G_S$-\Term{torsor}{挠子}时，$P$ 是否必然\Term{representable}{可表}？}

我们假定 $S$ 是\Term{connected}{连通}\Term{prescheme}{预概形}之和，即它的\Term{connected-component}{连通分支}都是开的；这立即把问题归结为 $S$ \Term{connected}{连通}的情形。此时选取 $S$ 的一个“\Term{geometric-point}{几何点}”$\xi$，即一个为\Term{algebraically-closed}{代数闭}域 $\Omega=\kappa(\xi)$ 的\Term{spectrum}{谱}的 $S$-\Term{scheme}{概形} $\xi$。对 $S$ 上群为 $G$ 的任意 \Term{galois-principal-bundle}{Galois 主纤维丛} $P$，$P_\xi$ 是\Term{algebraically-closed}{代数闭}域 $\kappa(\xi)$ 上的 \Term{galois-principal-bundle}{Galois 主纤维丛}，因此是\Term{trivial}{平凡}的。于是我们把最初的问题具体化为如下问题：
\SourcePage{88}
确定 $S$ 上在 $\xi$ 上方带点的 \Term{galois-principal-bundle}{Galois 主纤维丛}的\Term{category}{范畴}；也就是说，这些\Term{fibre-bundle}{纤维丛}配有一个 $S$-\Term{homomorphism}{同态} $\xi\to P$，即 $P_\xi$ 的一个\Term{trivialization}{平凡化}。固定 $G$ 后，把这类\Term{fibre-bundle}{纤维丛}在保持基点的\Term{isomorphism}{同构}下的类所成\Term{set}{集合}记作 $\pi^1(S,\xi;G)$。那么，$S$ 上群为 $G$ 的 \Term{galois-principal-bundle}{Galois 主纤维丛}（不指定基点）的\Term{isomorphism-class}{同构类}\Term{set}{集合} $\pi^1(S;G)$，\Term{isomorphism}{同构}于 $G$ 在 $\pi^1(S,\xi;G)$ 中的\Term{orbit}{轨道}\Term{set}{集合}；这里使用 $G$ 在此\Term{set}{集合}上的自然\Term{action}{作用}，它对应于用 $G$ \Term{translation}{平移}一个带点 \Term{galois-principal-bundle}{Galois 主纤维丛} $P$ 中的标记点：
\[
\pi^1(S;G)=\pi^1(S,\xi;G)/G.
\]
对任意\Term{morphism}{态射} $S'\to S$，假定它对于变动底上的\Term{twisted-constant-prescheme}{扭常值预概形}的\Term{fibered-category}{纤维范畴}是一个\Term{universal}{普遍}\Term{effective-descent}{有效下降}\Term{morphism}{态射}（例如 $S'\to S$ \Term{faithfully-flat}{忠实平坦}且\Term{locally-finite-presentation}{局部有限表示}，参见 5.4；其他例子参见 SGA 1 IX），我们要确定上述\Term{set}{集合}中分别记作 $\pi^1(S'/S,\xi;G)$ 和 $\pi^1(S'/S;G)$ 的\Term{subset}{子集}，它们由在 $S'$ 上成为\Term{trivial}{平凡}的（或称被 $S'$“\Term{split}{分裂}”的）$S$ 上 \Term{galois-principal-bundle}{Galois 主纤维丛}构成。实际上，我们将确定 $S$ 上被 $S'$ \Term{split}{分裂}的 \Term{galois-principal-bundle}{Galois 主纤维丛} $P$ 的\Term{category}{范畴}。当然，届时将有
\[
\pi^1(S,\xi;G)=\varinjlim_{S'}\pi^1(S'/S,\xi;G),
\]
其中右端的 $S'$ 遍历\Term{etale-topology}{平展拓扑}的\Term{covering-morphism}{覆盖态射} $S'/S$ 所成\Term{set}{集合}中的一个\Term{cofinal}{共尾}系统（例如当 $S$ \Term{quasi-compact}{拟紧}时，可取所有在 $S$ 上\Term{quasi-compact}{拟紧}、且结构\Term{morphism}{态射}\Term{etale}{平展}而\Term{surjective}{满射}的 $S'$）。同样，$S$ 上 \Term{galois-principal-bundle}{Galois 主纤维丛}的\Term{category}{范畴}，将是由 $S'$ 确定的子\Term{category}{范畴}（由在 $S'$ 上被\Term{split}{分裂}的\Term{fibre-bundle}{纤维丛}构成）的\Term{inductive-limit}{归纳极限}。

由关于 $S'/S$ 的假设，$S$ 上被 $S'$ \Term{split}{分裂}的 \Term{galois-principal-bundle}{Galois 主纤维丛}的\Term{category}{范畴}，\Term{equivalence}{等价}于 $S'$ 上配备相对于 $S'\to S$ 的\Term{descent-datum}{下降数据}的\Term{trivial}{平凡} \Term{galois-principal-bundle}{Galois 主纤维丛}的\Term{category}{范畴}（这些\Term{fibre-bundle}{纤维丛}形如 $G_{S'}$，其中 $G$ 以\Term{right-translation}{右平移}\Term{action}{作用}）。在被 $S'$ \Term{split}{分裂}的 $S$ 上 \Term{galois-principal-bundle}{Galois 主纤维丛} $P$ 上给定基点，用相应的 $S'$ 上\Term{trivial}{平凡}\Term{fibre-bundle}{纤维丛} $P'$ 及其\Term{descent-datum}{下降数据}来表述，就等于给定 $P'\times_{S'}S'_\xi$ 的一个\Term{trivialization}{平凡化}，且它须与相对于 $S'_\xi\to\xi$ 诱导的\Term{descent-datum}{下降数据}相容（注意这里置 $S'_\xi=S'\times_S\xi$）；即给定 $P'_\xi$ 在 $S'_\xi$ 上的一个与\Term{descent-datum}{下降数据}相容的\Term{section}{截面} $\sigma$。因此，对任意 $S$-\Term{prescheme}{预概形} $S'$（不再要求 $S'\to S$ 对\Term{twisted-constant-bundle}{扭常值纤维丛}的\Term{fibered-category}{纤维范畴}为\Term{universal}{普遍}\Term{effective-descent}{有效下降}\Term{morphism}{态射}……），把 $\pi^1(S'/S;G)$ 和 $\pi^1(S'/S,\xi;G)$ 定义为上述带\Term{descent-datum}{下降数据}的结构的\Term{isomorphism-class}{同构类}\Term{set}{集合}，是有益的。由此，对这些关于 $G$ 的\Term{functor}{函子}，我们得到一个极简单的\Term{simplicial}{单纯}描述，它用 $S'$ 在 $S$ 上的二重与三重\Term{fiber-product}{纤维积} $S''$ 和 $S'''$ 表述；我们将在下面勾勒这一描述（参见 6.3）。应当记住的重要结论如下。

\SourcePage{89}
\begin{SGAresult}{命题}{6.1}
假定 $S'$ 与 $S''$ 的\Term{connected-component}{连通分支}都是开的，并且，例如，$\pi_0(S')$ 对由两个\Term{projection}{投影} $\pi_0(S'')\to\pi_0(S')$ 诱导的\Term{equivalence-relation}{等价关系}所得\Term{quotient-set}{商集}只有一点。\footnote{编者注（48）：我们补加了 $S'$ 的\Term{connected-component}{连通分支}为开的假设以及第二个假设。后一个假设的意义是 6.3 所定义的\Term{simplicial-set}{单纯集} $K_\bullet=\pi_0(S_\bullet)$ \Term{connected}{连通}；实际上，只需更弱的假设：某个\Term{simplicial-set}{单纯集}\Term{morphism}{态射} $k_\bullet\to K_\bullet$ 的\Term{cone}{锥} $\widetilde K_\bullet$ \Term{connected}{连通}（参见该处）。}

(i) 从群\Term{category}{范畴}到\Term{set}{集合}\Term{category}{范畴}的\Term{functor}{函子} $G\mapsto\pi^1(S'/S,\xi;G)$ 由一个群\Term{representable}{表示}；记此群为 $\pi_1(S'/S,\xi)$，称为 $S$ 在 $\xi$ 处相对于 $S'\to S$ 的\Term{fundamental-group}{基本群}。因此有\Term{functor}{函子}性的\Term{bijective}{双射}
\[
\pi^1(S'/S,\xi;G)\simeq\operatorname{Hom}_{\mathrm{gr.}}(\pi_1(S'/S,\xi),G).
\]
(ii) 此群有一个与 $\pi_0(S'')$ \Term{bijective}{双射}的\Term{generator}{生成元}\Term{set}{集合}；用这些\Term{generator}{生成元}描述它的关系，则与 $\pi_0(S''')$ 的元素\Term{bijective}{双射}。\footnote{编者注（49）：我们删去了 $S'''$ 的\Term{connected-component}{连通分支}为开的多余假设。此外，下面给出的描述（参见\SGAsourcegap{原 PDF 第 89 页脚注（49）此处引注只有省略号，缺少具体指向}）以 $\pi_0(S'')\amalg\pi_0(S'_\xi)$ 为自然的\Term{generator}{生成元}\Term{set}{集合}；再利用来自二维胞腔的\Term{generator}{生成元}之间的关系，将其归结为 $\pi_0(S'')$。更精细的描述可参见例如 [Kan58]，第 19 节。} 特别地，若 $\pi_0(S'')$ \Term{finite}{有限}（分别地，若连同 $\pi_0(S''')$ 都\Term{finite}{有限}），则 $\pi^1(S'/S,\xi)$ \Term{finitely-generated}{有限生成}（分别地，\Term{finite-presentation}{有限表示}）。

(iii) $S$ 上被 $S'$ \Term{split}{分裂}、并在 $\xi$ 上方带基点的 \Term{galois-principal-bundle}{Galois 主纤维丛}的\Term{category}{范畴}，\Term{equivalence}{等价}于配备一个从 $\pi_1(S'/S,\xi)$ 到 $G$ 的\Term{homomorphism}{同态}的普通群 $G$ 的\Term{category}{范畴}。
\end{SGAresult}
证明将在下面给出，参见\SGAsourcegap{原 PDF 第 89 页命题 6.1 后“参见”之后只有省略号，缺少具体引注}。

\textbf{6.2。} 当 $S$ 是\Term{connected}{连通}\Term{locally-noetherian}{局部诺特}\Term{prescheme}{预概形}时，$S$ 上任意\Term{etale}{平展}\Term{prescheme}{预概形} $S'$ 都是\Term{locally-noetherian}{局部诺特}的，所以其\Term{connected-component}{连通分支}都是开的。由上述结果可\Term{pushout}{推出}\footnote{编者注（50）：可以详细说明这一步推论；$I$ 与\Term{etale-topology}{平展拓扑}的\Term{covering-morphism}{覆盖态射} $S'\to S$ 的一个\Term{cofinal}{共尾}\Term{set}{集合}\Term{bijective}{双射}……}，从普通群\Term{category}{范畴}到\Term{set}{集合}\Term{category}{范畴}的\Term{functor}{函子} $G\mapsto\pi^1(S,\xi;G)$ 是\Term{strict-prorepresentable}{严格前可表}的（参见《Bourbaki 研讨班》，1960 年 2 月，第 195 号，A.2 与 A.3 节）；即在一个向上\Term{filtered-family}{滤过}的指标集 $I$ 上，存在一个普通群的\Term{projective-system}{投射系}
\[
\Pi=\Pi_1(S;\xi)=(\pi_i)_{i\in I},
\]
它是“严格”的（即转移\Term{morphism}{态射} $\pi_j\to\pi_i$ 为\Term{surjective}{满射}），并且有关于 $G$ 的\Term{functor}{函子}\Term{isomorphism}{同构}
\[
\pi^1(S,\xi;G)\simeq\varinjlim_i\operatorname{Hom}_{\mathrm{gr.}}(\pi_i,G).
\]
右端也简记为 $\operatorname{Hom}_{\mathrm{pro\text{-}gr.}}(\Pi,G)$（参见该处）。

在\Term{projective-limit}{投射极限} $\pi=\varprojlim\pi_i$“\Term{large-relative-to-cardinal}{足够大}”的情形，确切地说，当\Term{canonical}{典范}\Term{homomorphism}{同态} $\Pi\to\Pi_i$ 都为\Term{surjective}{满射}时，应给
\SourcePage{90}
$\Pi$ 配备各 $\Pi_i$ 的\Term{discrete}{离散}\Term{topology}{拓扑}的\Term{projective-limit}{投射极限}\Term{topology}{拓扑}，而上述\Term{isomorphism}{同构}也可写成
\[
\pi^1(S,\xi;G)\simeq\operatorname{Hom}_{\mathrm{gr.\ top.}}(\pi,G),
\]
其中右端表示\Term{topological-group}{拓扑群}\Term{homomorphism}{同态}的\Term{set}{集合}，约定给 $G$ 配备\Term{discrete}{离散}\Term{topology}{拓扑}。

众所周知，当 $\pi_i$ 都是\Term{finite}{有限}群时，我们刚对\Term{projective-system}{投射系} $\Pi$ 提出的假设成立（参见 [BEns]，III，第 7.4 节，\Term{theorem}{定理} 1）。显然，这最后一个条件也就是说，$S$ 上每个 \Term{galois-principal-bundle}{Galois 主纤维丛}都是\Term{isotrivial}{等平凡}的，即被一个\Term{finite}{有限}的\Term{etale}{平展}\Term{surjective}{满射}\Term{morphism}{态射}\Term{split}{分裂}。当 $S$ \Term{geometrically-unibranch}{几何单分支}（例如\Term{normal}{正规}）时就是如此，这由 5.14 立即得出。\footnote{编者注（51）：已将 5.13 改为 5.14。} 在 $\pi_i$ 都\Term{finite}{有限}时，群 $\pi$ 也与 SGA 1 V 中引进的\Term{fundamental-group}{基本群} $\pi_1(S,\xi)$ 相同。

因此，在有利的情形（$\pi\to\pi_i$ 都为\Term{surjective}{满射}），可称 $\pi$ 为 $S$ 在 $\xi$ 处的\Term{enlarged-fundamental-group}{扩大基本群}。在此有利情形之外，$\pi$ 本身几乎没有意义，通常的\Term{fundamental-group}{基本群}的角色由\Term{projective-system}{投射系} $\Pi$ 本身担当；我们称它为 $S$ 在 $\xi$ 处的\Term{enlarged-fundamental-progroup}{扩大基本前群}。（欢迎提出比“扩大”更好的术语！\footnote{编者注（52）：有些作者称之为“真正的”\Term{fundamental-group}{基本群}。}）须指出，知道这个\Term{progroup}{前群}，比知道 SGA 1 V 的通常\Term{fundamental-group}{基本群} $\pi_1(S,\xi)$ 更为精细；确切地说，后者是各 $\pi_i$ 的\Term{finite}{有限}商组成的\Term{projective-system}{投射系}的\Term{projective-limit}{投射极限}。

\textbf{6.3。} 简要说明 $\pi_1(S'/S,\xi)$ 的“计算”。设 $S_i$ 是 $S'$ 在 $S$ 上的第 $(i+1)$ 次\Term{fiber-product}{纤维幂}（即 $S_0=S'$、$S_1=S''$ 等）。各 $S_i$ 之间有显然的\Term{simplicial}{单纯}运算，使 $(S_i)_{i\in\mathbb N}$ 成为 $(\mathrm{Sch})_{/S}$ 中的一个\Term{simplicial-object}{单纯对象}。
把这个\Term{simplicial-object}{单纯对象}用“\Term{connected-component}{连通分支}\Term{set}{集合}”\Term{functor}{函子}
\[
\pi_0:(\mathrm{Sch})_{/S}\longrightarrow(\mathrm{Ens})
\]
变换，便得到\Term{simplicial-set}{单纯集} $K_\bullet=(K_i)_{i\in\mathbb N}$，其中 $K_i=\pi_0(S_i)$。
同样，$S_{i\xi}$（即 $S'_\xi$ 在 $\xi$ 上的第 $(i+1)$ 次\Term{fiber-product}{纤维幂}）构成 $(\mathrm{Sch})_{/\xi}$，因而也构成 $(\mathrm{Sch})_{/S}$ 中的一个\Term{simplicial-object}{单纯对象}；它还带有到 $(S_i)_{i\in\mathbb N}$ 的自然\Term{simplicial-object}{单纯对象}\Term{homomorphism}{同态}。由此得到一个\Term{simplicial-set}{单纯集} $k_\bullet$（其中 $k_i=\pi_0(S_{i\xi})$）和一个\Term{canonical}{典范}\Term{homomorphism}{同态}
\[
k_\bullet\longrightarrow K_\bullet.
\]
取此\Term{morphism}{态射}的\Term{cone}{锥}，可构成一个新的\Term{simplicial-set}{单纯集}（参见 9.5.1）：
\[
\widetilde K_\bullet=\operatorname{Cone}(k_\bullet\longrightarrow K_\bullet).
\]
\SourcePage{91}
\footnote{编者注（53）：原稿考虑 $\widetilde K_\bullet=K_\bullet/k_\bullet$，我们作了修正；一方面，$k_\bullet$ 本来就\Term{contractible}{可缩}（参见 9.7.1）；另一方面，\Term{morphism}{态射} $k_\bullet\to K_\bullet$ 一般并不\Term{injectivity}{单射}。当 $S'/S$ \Term{finite-etale}{有限平展}时，它是\Term{epimorphism}{满态射}（参见 9.7.3）。} 这样就得到一个“\Term{pointed-simplicial-set}{带点单纯集}”$\widetilde K_\bullet$（即带有\Term{homomorphism}{同态} $\widetilde\xi:e_\bullet\to\widetilde K_\bullet$ 的\Term{simplicial-set}{单纯集}，其中 $e_\bullet$ 是终\Term{simplicial-set}{单纯集}）。可以构造其众所周知的组合\Term{invariants}{不变量} $\pi_0(\widetilde K_\bullet,\widetilde\xi)$ 和 $\pi_1(\widetilde K_\bullet,\widetilde\xi)$；其构造实际上只涉及\Term{degree}{次数} $\leq1$ 的部分，分别地，\Term{degree}{次数} $\leq2$ 的部分。定义这些\Term{invariants}{不变量}，无须对 $S$ 或 $S'$ 施加\Term{restriction}{限制}。当 $S_0$ 与 $S_1$ 的\Term{connected-component}{连通分支}都开，且 $\widetilde K_\bullet$ \Term{connected}{连通}时，\footnote{编者注（54）：我们补加了 $\widetilde K_\bullet$ \Term{connected}{连通}这一假设；证明见下面的增补（第 9 节）。} 不难验证 $\pi_1(\widetilde K_\bullet,\widetilde\xi)$ \Term{representable}{表示}\Term{functor}{函子} $G\mapsto\pi^1(S'/S,\xi;G)$，即有
\[
\pi_1(S'/S;\xi)\simeq\pi_1(\widetilde K_\bullet,\widetilde\xi).
\]
还应指出，当\Term{morphism}{态射} $S'\to S$“\Term{universally-submersive}{普遍淹没}”（参见 SGA 1，IV 2.1），且 $S'$ 的\Term{connected-component}{连通分支}都开时，\footnote{编者注（55）：我们根据上面所作的修正改写了后面的文字，参见编者注（53）。（原文是：“$\pi_0(\widetilde K_\bullet,\widetilde\xi)$ 也\Term{canonical}{典范}\Term{isomorphism}{同构}于 $S$ 的\Term{connected-component}{连通分支}\Term{set}{集合} $\pi_0(S,\xi)$，以 $\xi$ 在 $S$ 中所在的\Term{connected-component}{连通分支}为标记点。”）} \Term{simplicial-set}{单纯集} $K_\bullet$，因而 $\widetilde K_\bullet$ 也是\Term{connected}{连通}的。

\textbf{例 6.4。} 还需给出\Term{enlarged-fundamental-group}{扩大基本群}的例子。回到 1.6 的例子：设 $k$ 是\Term{algebraically-closed}{代数闭}域，$C_1$ 是 $k$ 上\Term{complete}{完备}有理\Term{algebraic-curve}{曲线}，恰有一个\Term{singularity}{奇点}，且它是\Term{ordinary-double-point}{普通二重点}；$C_2$ 是两个\Term{irreducible-component}{不可约分支}之并，这两个分支均\Term{isomorphism}{同构}于 $\mathbb P^1_k$，且恰好相交于两点，这两点都是 $C_2$ 的\Term{ordinary-double-point}{普通二重点}。在这两种情形，\Term{curve}{曲线}的\Term{enlarged-fundamental-group}{扩大基本群}都\Term{isomorphism}{同构}于 $\mathbb Z$。\footnote{编者注（56）：可以详细说明这一点：首先，由 5.14，\Term{projective-line}{射影直线} $\mathbb P^1_k$ 的\Term{enlarged-fundamental-group}{扩大基本群}为零，即 $\mathbb P^1_k$ 是“真正”\Term{simply-connected}{单连通}的。其次，需要把 SGA 1 IX 的推论 5.4 及其后的讨论扩展到\Term{enlarged-fundamental-group}{扩大基本群}与\Term{principal-bundle}{主纤维丛}（不必\Term{finite}{有限}）的\Term{category}{范畴}的情形。（设 $\Gamma_1$ 和 $\Gamma_2$ 是 $\mathbb P_k$ 的两个副本，$a_i,b_i$ 是 $\Gamma_i$ 的两个不同点，$C''_2$ 是 $\Gamma_1$ 与 $\Gamma_2$ 的\Term{disjoint-union}{不交并}，$C'_2$ 是将 $a_1$ 与 $a_2$ \Term{gluing}{粘合}得到的\Term{curve}{曲线}；则再将 $b_1$ 与 $b_2$ \Term{gluing}{粘合}就从 $C'_2$ 得到 $C_2$。把该处之后的讨论扩展到扩大情形，便表明 $C'_2$（分别地 $C_2$）的\Term{enlarged-fundamental-group}{扩大基本群}为零（分别地为 $\mathbb Z$）。}

一般而言，应在\Term{enlarged-fundamental-group}{扩大基本群}的框架中，重新处理 SGA 1 IX 5 的结果（并加以简化和修正）。该处 5.5 的例子会给出同样多的并非\Term{profinite}{前有限}的\Term{enlarged-fundamental-progroup}{扩大基本前群}的例子。例如，在第一个例子中，若把\Term{ordinary-double-point}{普通二重点}换成具有 $n$ 个不同分支的二重点，\footnote{编者注（57）：这里指具有 $n$ 条不同切线的 $n$ 重点，例如\Term{curve}{曲线} $X^n-Y^n=X^{n+1}$。} 就会得到有 $n-1$ 个\Term{generator}{生成元}的\Term{free-group}{自由群}（\Term{discrete}{离散}的！）作为\Term{enlarged-fundamental-group}{扩大基本群}。

\SourcePage{92}
\SGAsection{7}{用扩大基本群分类扭常值预概形及有限型乘法型群}
\textbf{7.0。} 设 $S$ 是一个\Term{prescheme}{预概形}，仍假定它\Term{locally-noetherian}{局部诺特}，以保证 $S$ 以及我们将考虑的 $S$ 上某些\Term{prescheme}{预概形}（尤其是 $S$ 上的\Term{etale}{平展}\Term{prescheme}{预概形}，更一般地，$S$ 上\Term{locally-finite-type}{局部有限型}的\Term{prescheme}{预概形}）都是\Term{locally-connected}{局部连通}的。

\begin{SGAresult}{命题}{7.0.1}
$S$ 上每个对于（fppf）\Term{topology}{拓扑}\Term{locally-trivial}{局部平凡}的\Term{twisted-constant-prescheme}{扭常值预概形} $X$（即被一个\Term{faithfully-flat}{忠实平坦}\Term{locally-finite-presentation}{局部有限表示}\Term{morphism}{态射} $S'\to S$ \Term{split}{分裂}），都是\Term{quasi-isotrivial}{拟等平凡}的（即甚至可选 $S'\to S$ 为\Term{etale}{平展}\Term{surjective}{满射}）。
\end{SGAresult}
\begin{proof}
可以假定 $S$ \Term{connected}{连通}，因而 $X$ 为型 $I$，其中 $I$ 是一个固定\Term{set}{集合}。于是 $X'=X\times_SS'$ \Term{isomorphism}{同构}于 $I_{S'}$，所以 $I_{S'}$ 带有相对于 $S'\to S$ 的一个\Term{descent-datum}{下降数据}，即一个满足通常\Term{transitivity}{传递性}条件的\Term{isomorphism}{同构} $I_{S''}\xrightarrow\sim I_{S''}$。而 $S''=S'\times_SS'$ \Term{locally-noetherian}{局部诺特}，所以\Term{locally-connected}{局部连通}；因此，$I_{S''}$ 的\Term{automorphism}{自同构}对应于 $G_{S''}$ 的\Term{section}{截面}，其中 $G=\operatorname{Aut}(I)$ 是 $I$ 的\Term{permutation-group}{置换群}。
这样就在 $G_{S'}$（视为\Term{trivial}{平凡} \Term{galois-principal-bundle}{Galois 主纤维丛}）上得到相对于 $S'\to S$ 的一个\Term{descent-datum}{下降数据}。由 5.4，此\Term{descent-datum}{下降数据}是\Term{effective}{有效}的，从而给出 $S$ 上群为 $G$ 的一个 \Term{galois-principal-bundle}{Galois 主纤维丛} $P$。按构造，在 $S$ 上的\Term{locally-noetherian}{局部诺特}\Term{prescheme}{预概形}\Term{category}{范畴}中，它\Term{representable}{表示}\Term{functor}{函子} $\operatorname{Isom}_S(I_S,X)$。因此，\Term{etale}{平展}\Term{surjective}{满射}\Term{base-change}{基变换} $P\to S$ \Term{split}{分裂} $X$，所以 $X$ 确为\Term{quasi-isotrivial}{拟等平凡}。
\end{proof}
\textbf{注记 7.0.2。} 应当注意，即使 $S$ 是一个域的\Term{spectrum}{谱}，也不能说 $S$ 上每个\Term{twisted-constant-bundle}{扭常值纤维丛}都是\Term{quasi-isotrivial}{拟等平凡}的。例如，只需取 $X$ 为一列形如 $\operatorname{Spec}(k_i)$ 的\Term{scheme}{概形}之和，其中 $k_i$ 是 $k$ 的\Term{degree}{次数}严格递增的\Term{separable-extension}{可分扩张}。

上面的证明同时表明，$S$ 上型为 $I$ 的\Term{quasi-isotrivial}{拟等平凡}\Term{twisted-constant-bundle}{扭常值纤维丛} $X$ 的分类，\Term{equivalence}{等价}于 $S$ 上群为 $G=\operatorname{Aut}(I)$ 的 \Term{galois-principal-bundle}{Galois 主纤维丛}的分类。这甚至是一个\Term{category}{范畴}\Term{equivalence}{等价}。

如 SGA 1 V 中一样，可把它写成更方便的\Term{form}{形式}。为此，假定 $S$ \Term{connected}{连通}且带有一个\Term{geometric-point}{几何点} $\xi$。于是\Term{enlarged-fundamental-progroup}{扩大基本前群} $\Pi=\Pi_1(S,\xi)$ 有定义。另一方面，对 $S$ 上任意\Term{quasi-isotrivial}{拟等平凡}\Term{twisted-constant-bundle}{扭常值纤维丛} $X$，设 $I=X(\xi)$ 是它在 $\xi$ 处的\Term{set-theory}{集合论}\Term{fiber}{纤维}。则 $X$ 为型 $I$，从而如刚才所述，它与 $S$ 上群为 $G=\operatorname{Aut}(I)$ 的 \Term{galois-principal-bundle}{Galois 主纤维丛} $P=\operatorname{Isom}_S(I_S,X)$ 相联系。
根据 $\Pi$ 的定义，得到一个从 $\Pi$ 到 $G$ 的\Term{canonical}{典范}\Term{homomorphism}{同态}，即某个 $\Pi_i$ 到 $G$ 的\Term{homomorphism}{同态}。由于 $G$ 是 $I=X(\xi)$ 的\Term{permutation-group}{置换群}，这意味着 $\Pi$“在 $I=X(\xi)$ 上\Term{continuous}{连续}\Term{action}{作用}”；这里的含义是，对充分大的 $i$，$\pi_i$ 在 $I$ 上\Term{action}{作用}，并与转移\Term{morphism}{态射}相容。
\SourcePage{93}
我们留待读者验证，$S$ 上\Term{quasi-isotrivial}{拟等平凡}\Term{twisted-constant-bundle}{扭常值纤维丛}之间的每个 $S$-\Term{morphism}{态射} $X\to Y$，都诱导一个与 $\Pi$ 的\Term{action}{作用}相容的映射 $X(\xi)\to Y(\xi)$，而这样得到的\Term{functor}{函子}是一个\Term{category}{范畴}\Term{equivalence}{等价}：
\begin{SGAresult}{命题}{7.0.3}
设 $S$ 是\Term{locally-noetherian}{局部诺特}\Term{connected}{连通}\Term{prescheme}{预概形}，$\xi$ 是 $S$ 的一个\Term{geometric-point}{几何点}，$\Pi=\Pi_1(S,\xi)$ 是 $S$ 在 $\xi$ 处的\Term{enlarged-fundamental-progroup}{扩大基本前群}。则\Term{functor}{函子}
\[X\longmapsto X(\xi)\]
给出 $S$ 上\Term{quasi-isotrivial}{拟等平凡}\Term{twisted-constant-bundle}{扭常值纤维丛}的\Term{category}{范畴}与 $\Pi$ 在其上\Term{continuous}{连续}\Term{action}{作用}的\Term{set}{集合}的\Term{category}{范畴}之间的\Term{equivalence}{等价}。
\end{SGAresult}
此\Term{functor}{函子}与\Term{finite}{有限}和及\Term{finite}{有限}\Term{projective-limit}{投射极限}运算相容。因此，例如，$S$ 上\Term{quasi-isotrivial}{拟等平凡}\Term{twisted-constant-group}{扭常值群}（或环等），对应于\Term{progroup}{前群} $\Pi$ 在其上\Term{continuous}{连续}\Term{action}{作用}的普通群（分别为环等）。特别地：
\begin{SGAresult}{推论}{7.0.4}
$S$ 上\Term{quasi-isotrivial}{拟等平凡}\Term{twisted-constant-group}{扭常值}\Term{commutative-group}{交换群}的\Term{category}{范畴}，\Term{equivalence}{等价}于“$\Pi$-\Term{module}{模}”的\Term{category}{范畴}，即 $\Pi$ 在其上\Term{continuous}{连续}\Term{action}{作用}的\Term{commutative-group}{交换群} $M$ 的\Term{category}{范畴}。
\end{SGAresult}
现在使用 5.7，得到：
\begin{SGAresult}{\Term{theorem}{定理}}{7.1}
设 $S$ 是\Term{locally-noetherian}{局部诺特}\Term{connected}{连通}\Term{prescheme}{预概形}，$\xi$ 是 $S$ 的一个\Term{geometric-point}{几何点}，$\Pi=\Pi_1(S,\xi)$ 是 $S$ 在 $\xi$ 处的\Term{enlarged-fundamental-progroup}{扩大基本前群}。则\Term{functor}{函子}
\[
G\longmapsto\operatorname{Hom}_{\kappa(\xi)\text{-}\mathrm{gr.}}(G_\xi,\mathbb G_{m,\xi})
\]
诱导 $S$ 上\Term{quasi-isotrivial}{拟等平凡}\Term{multiplicative-type}{乘法型群}的\Term{category}{范畴}与 $\Pi$-\Term{module}{模}的\Term{category}{范畴}之间的\Term{anti-equivalence}{反等价}。
\end{SGAresult}
使用 4.5，得到：
\begin{SGAresult}{推论}{7.2}
上述\Term{functor}{函子}诱导 $S$ 上\Term{finite-type}{有限型}\Term{multiplicative-type}{乘法型群}的\Term{category}{范畴}与在 $\mathbb Z$ 上\Term{finitely-generated}{有限生成}的 $\Pi$-\Term{module}{模}的\Term{category}{范畴}之间的\Term{anti-equivalence}{反等价}。
\end{SGAresult}
\textbf{例 7.3。} 例如，再取 $S$ 为\Term{algebraically-closed}{代数闭}域上的一条\Term{complete}{完备}有理\Term{curve}{曲线}，恰有一个具有 $n+1$ 个不同分支的重点。由 6.4，\Term{enlarged-fundamental-group}{扩大基本群} $\Pi(S,\xi)$ 是有 $n$ 个\Term{generator}{生成元}的\Term{free-group}{自由群}。因此，由 7.2，$S$ 上相对\Term{dimension}{维数}为 $m$ 的\Term{torus}{环面}的分类，\Term{equivalence}{等价}于 $\mathbb Z$-\Term{module}{模} $\mathbb Z^m$ 的 $n$ 个\Term{endomorphism}{自同态} $A_1,\ldots,A_n$ 所成系统在 $\mathbb Z^m$ 的\Term{automorphism}{自同构}下的分类。除 $n\leq1$ 或 $m\leq1$ 外，明确分类这些系统似乎毫无希望。至少能为这样的系统定义大量非\Term{trivial}{平凡}\Term{invariants}{不变量}，例如各 $A_i$ 的\Term{characteristic-polynomial}{特征多项式}。\footnote{编者注（58）：特别地，这表明 1.6 所考虑的相对\Term{dimension}{维数}为 $2$ 的\Term{torus}{环面}并非\Term{isotrivial}{等平凡}。}
\SourcePage{94}
\textbf{注记 7.4。} 若不对 $S$ 作任何假设，仍有如下事实：给定一个在 $\mathbb Z$ 上\Term{finitely-generated}{有限生成}的普通\Term{commutative-group}{交换群} $M$，$S$ 上型为 $M$ 的\Term{multiplicative-type}{乘法型群}的\Term{category}{范畴}，\Term{anti-equivalence}{反等价}于 $S$ 上群为 $G=\operatorname{Aut}_{\mathrm{gr.}}(M)$ 的 \Term{galois-principal-bundle}{Galois 主纤维丛}的\Term{category}{范畴}。这容易由 5.9 和 5.10 得出。

\textbf{注记 7.5。} 我们在这两节中勾勒的\Term{fundamental-progroup}{基本前群}理论，最好在一般\Term{site}{景}的框架下写出。以此\Term{form}{形式}，它也适用于例如通常的\Term{topological-space}{拓扑空间}，而且至少对\Term{locally-connected}{局部连通}的\Term{topological-space}{拓扑空间}（不必\Term{local-scheme}{局部}\Term{simply-connected}{单连通}），给出令人满意的理论。在此情形，似乎也不能满足于只定义一个\Term{fundamental-group}{基本群}，而需要一个\Term{progroup}{前群}。最后指出，一旦具备\Term{topology}{拓扑}与\Term{descent}{下降}的语言（这才是这些问题真正的根本），这里勾勒的论述在技术上也比 SGA 1 V 更简单，因此原则上应当替代后者。

\SGAsection{8}{附录：去掉某些仿射性假设}
我们的目的是证明 4.9 的如下推广。
\begin{SGAresult}{\Term{theorem}{定理}}{8.1}
设 $S$ 是一个\Term{prescheme}{预概形}，$H$ 是 $S$ 上\Term{flat}{平坦}\Term{finite-presentation}{有限表示}的\Term{group-prescheme}{群预概形}，其\Term{fiber}{纤维}\Term{connected}{连通}且\Term{affine}{仿射}。\footnote{编者注（59）：注意，由第 VI\textsubscript{B} 讲\Term{theorem}{定理} 5.3 或推论 5.5，这些假设蕴涵 $H$ 在 $S$ 上\Term{separated}{分离}。} 设 $s\in S$，假定 $H_s$ 是\Term{torus}{环面}。则存在 $s$ 的一个\Term{open-neighborhood}{开邻域} $U$，使 $H|U$ 是\Term{torus}{环面}。
\end{SGAresult}
立即得到：
\begin{SGAresult}{推论}{8.2}
设 $H$ 是 $S$ 上\Term{flat}{平坦}\Term{finite-presentation}{有限表示}的\Term{group-prescheme}{群预概形}。$H$ 是\Term{torus}{环面}的充要条件是其\Term{fiber}{纤维}都是\Term{torus}{环面}。
\end{SGAresult}
\textbf{注记 8.3。} 即使 $S$ 是一个\Term{discrete-valuation-ring}{离散赋值环}的\Term{spectrum}{谱}，在 8.1 中也不能去掉 $H$ 的\Term{fiber}{纤维}（在此即\Term{generic-fiber}{泛纤维}）\Term{affine}{仿射}的假设，因为存在 $S$ 上\Term{smooth}{光滑}群的例子，其\Term{generic-fiber}{泛纤维}是\Term{elliptic-curve}{椭圆曲线}，而\Term{special-fiber}{特殊纤维}是 $\mathbb G_m$。

\begin{proof}[8.1 的证明]
显然可假定 $S=\operatorname{Spec}(A)$ \Term{affine}{仿射}；用标准方法（参见 EGA IV\textsubscript{2}，8.8.2），便归结为 $S$ 还\Term{noetherian}{诺特}的情形。我们先在此情形证明 8.2。
由 4.9，只需证明 $H$ 在 $S$ 上\Term{affine}{仿射}。于是可假定\footnote{编者注（60）：再用一次 EGA IV\textsubscript{2}，8.8.2。} $A$ 是\Term{local-noetherian-ring}{局部诺特环}；又因为\Term{completion}{完备化} $\widehat A$ 在 $A$ 上\Term{faithfully-flat}{忠实平坦}，由\Term{descent}{下降}可归结为 $A$ 是\Term{complete}{完备}\Term{local-noetherian-ring}{局部诺特环}的情形。设 $S'$ 是 $S_{\mathrm{r\text{é}d}}$ 的\Term{normalization}{正规化}；由 Nagata 的结果，知道它在 $S$ 上\Term{finite}{有限}（EGA $0_{\mathrm{IV}}$，23.1.5）；此外 $S'\to S$ 为\Term{surjective}{满射}，所以 $H'=H\times_SS'\to H$ \Term{finite}{有限}且\Term{surjective}{满射}；于是要证明 $H$ \Term{affine}{仿射}，只需证明 $H'$ 如此（EGA II，6.7.1）。
\SourcePage{95}
\footnote{编者注（61）：在后文中，我们详细展开了原稿。} 把 $S$ 换成 $S'$ 的一个\Term{connected-component}{连通分支}，便归结为 $A$ 是\Term{normal}{正规}且\Term{integral-ring}{整}的\Term{noetherian-ring}{诺特环}（\Term{semilocal-ring}{半局部环}）的情形。再把 $A$ 换成它在其\Term{fraction-field}{分式域}的一个\Term{finite-separable-extension}{有限可分扩张}中的\Term{normalization}{正规化}，可假定 $H$ 的\Term{generic-fiber}{泛纤维} $H_\eta$ \Term{diagonalizable}{可对角化}，即有\Term{isomorphism}{同构}
\[u_\eta:H_\eta\xrightarrow\sim T_\eta,\qquad T=\mathbb G_{m,S}^r.\]
此时有如下\Term{lemma}{引理}。
\end{proof}
\begin{SGAresult}{\Term{lemma}{引理}}{8.4}
设 $S$ 是\Term{locally-noetherian}{局部诺特}\Term{normal}{正规}\Term{irreducible}{不可约}\Term{prescheme}{预概形}，\Term{generic-point}{泛点}为 $\eta$；$H$ 是 $S$ 上\Term{smooth}{光滑}且\Term{fiber}{纤维}\Term{connected}{连通}的\Term{group-prescheme}{群预概形}，$T$ 是 $S$ 上\Term{finite-type}{有限型}\Term{multiplicative-type}{乘法型}\Term{group-prescheme}{群预概形}，$u_\eta:H_\eta\to T_\eta$ 是 $\kappa(\eta)$ 上的\Term{algebraic-group}{代数群}\Term{homomorphism}{同态}。则 $u_\eta$ 延拓为一个群\Term{homomorphism}{同态} $u:H\to T$。
\end{SGAresult}
\begin{proof}
把 $S$ 换成它在其\Term{function-field}{函数域}的一个\Term{finite-separable-extension}{有限可分扩张}中的\Term{normalization}{正规化}，可假定 $T$ \Term{diagonalizable}{可对角化}（我们所要作的应用本就如此）。于是 $T$ 是形如 $\mathbb G_{m,S}^r$ 的群的一个\Term{closed-subgroup}{闭子群}，从而归结为 $T=\mathbb G_{m,S}$ 的情形。
问题全在于证明，将 $u_\eta$ 视为从 $H$ 到 $T=\mathbb G_{m,S}$ 的\Term{rational-map}{有理映射}时，它处处有定义（因为延拓它的\Term{morphism}{态射} $u:H\to T$ 届时必然为群\Term{homomorphism}{同态}）。可把 $u_\eta$ 看作 $H_\eta$ 的\Term{structure-sheaf}{结构层}的一个\Term{invertible}{可逆}\Term{section}{截面} $f$，而需证明它可延拓为 $H$ 的\Term{structure-sheaf}{结构层}的一个\Term{invertible}{可逆}\Term{section}{截面}。$H$ 在\Term{normal}{正规}的 $S$ 上\Term{smooth}{光滑}，因而\Term{normal}{正规}（SGA 1，II 3.1）；所以只需找到 $H$ 的一个\Term{codimension}{余维} $\geq2$ 的\Term{closed-subset}{闭子集} $Z$，使 $f$ 可延拓为 $H-Z$ 的\Term{structure-sheaf}{结构层}的\Term{invertible}{可逆}\Term{section}{截面}。这立即把问题归结为 $S$ 是\Term{discrete-valuation-ring}{离散赋值环} $A$ 的\Term{spectrum}{谱}的情形（在 $S$ 的\Term{codimension}{余维} $1$ 的点处\Term{localization}{局部化}）。

设 $t$ 是 $A$ 的一个\Term{uniformizer}{一致化参数}，$t'$ 是它所定义的 $\mathcal O_H$ 的\Term{section}{截面}，故\Term{special-fiber}{特殊纤维} $H_0$ 等于 $V(t')$。按假设，$H_0$ 在\Term{residue-field}{剩余域} $k$ 上\Term{smooth}{光滑}且\Term{connected}{连通}。于是 $f$ 是 $H$ 上的一个\Term{rational-function}{有理函数}，在 $H-H_0$ 上既无零点又无极点；由于 $H_0=V(t')$ 是\Term{irreducible}{不可约}\Term{divisor}{除子}，存在整数 $n\in\mathbb Z$，使 $t'^nf=t^nf$ 既无零点又无极点，即为 $\mathcal O_H$ 的一个\Term{invertible}{可逆}\Term{section}{截面}。因此它定义一个\Term{morphism}{态射} $v:H\to\mathbb G_{m,S}$；又因 $v_\eta=t'^nu_\eta$ 且 $u_\eta$ 为群\Term{homomorphism}{同态}，$v$ 把 $H$ 的\Term{unit-section}{单位截面}变成 $\mathbb G_{m,S}$ 的一个\Term{section}{截面}，其在 $S$ 的\Term{generic-point}{泛点}处的值为 $t^n$；由于这是 $\mathbb G_{m,S}$ 的\Term{section}{截面}，$t^n$ 必须为单位，即 $n=0$。所以 $v$ 延拓 $u_\eta$，这就证明了 8.4。
\end{proof}
将此\Term{lemma}{引理}用于当前情形，得到一个群\Term{homomorphism}{同态}
\[u:H\longrightarrow T=\mathbb G_{m,S}^r,\]
它在\Term{generic-fiber}{泛纤维}上诱导\Term{isomorphism}{同构}。我们证明 $u$ 是\Term{isomorphism}{同构}。

\SourcePage{96}
\begin{SGAresult}{\Term{lemma}{引理}}{8.5}
对每个与 $A$ 的\Term{residue-characteristic}{剩余特征}互素的整数 $n>0$，${}_nH=\operatorname{Ker}(n\cdot\operatorname{id}_H)$ 在 $S$ 上\Term{finite}{有限}。
\end{SGAresult}
\begin{proof}
若 $n$ 与 $A$ 的\Term{residue-characteristic}{剩余特征}互素，它便与 $S$ 各点的所有\Term{residue-characteristic}{剩余特征}互素。所以 $n\cdot\operatorname{id}_H$ 在 $H$ 的每个\Term{fiber}{纤维}上诱导\Term{etale}{平展}\Term{morphism}{态射}，因而 $n\cdot\operatorname{id}_H$ \Term{etale}{平展}（SGA 1，I 5.9），所以其\Term{kernel}{核} ${}_nH$ 在 $S$ 上\Term{etale}{平展}。另一方面，${}_nH$ 在 $S$ 上\Term{separated}{分离}，因为 $H$ 如此。\footnote{原注（$*$）：由 Raynaud \Term{theorem}{定理} VI\textsubscript{B} 5.3。}\footnote{编者注（62）：结合 VI\textsubscript{B} 第 5 节的修改，详细说明这一点。} 此外，其所有\Term{fiber}{纤维}的\Term{rank}{秩}相同，均为 $n^r$，因为 $H$ 的\Term{fiber}{纤维}均为\Term{torus}{环面}，且\Term{dimension}{维数}同为 $r$（$H$ 在 $S$ 上\Term{smooth}{光滑}）。由此\Term{pushout}{推出} ${}_nH$ 在 $S$ 上\Term{finite}{有限}（SGA 1，I 10.9 或 EGA IV\textsubscript{4}，18.2.9）。
\end{proof}
所以，由上面的\Term{lemma}{引理}，$u({}_nH)$ 是 $T$ 的一个\Term{closed-subset}{闭子集}。由于此\Term{set}{集合}在\Term{generic-fiber}{泛纤维} $T_\eta$ 上等于 ${}_n(T_\eta)$，它包含后者的\Term{closure}{闭包}，即 ${}_nT$。对任意 $s\in S$，各 ${}_n(T_s)$ 在 $T_s$ 中\Term{dense}{稠密}；而 $u_s(H_s)$ 是包含它们的一个\Term{closed-subset}{闭子集}（VI\textsubscript{B} 1.4.2），所以 $u_s(H_s)=T_s$，因而 $u$ 为\Term{surjective}{满射}。域上两个同\Term{dimension}{维数}\Term{torus}{环面}之间的\Term{surjective}{满射}\Term{homomorphism}{同态}是\Term{flat}{平坦}的，\footnote{编者注（63）：这例如由 VIII 3.2 a) 得到；更一般地，若 $f:G\to H$ 是域 $k$ 上的\Term{algebraic-group}{代数群}的\Term{surjective}{满射}\Term{morphism}{态射}，且 $H$ \Term{reduced}{约化}，则 $f$ \Term{flat}{平坦}（参见 VI\textsubscript{B} 1.3）。} 因而 $u$ 在每个\Term{fiber}{纤维}上诱导\Term{flat}{平坦}\Term{morphism}{态射}，所以 $u$ \Term{flat}{平坦}（SGA 1 I 5.9）。于是 $K=\operatorname{Ker}(u)$ 在 $S$ 上\Term{flat}{平坦}，\footnote{编者注（64）：在后文中，我们详细展开了原稿。} 因而等于其\Term{generic-fiber}{泛纤维} $K_\eta$ 的\Term{closure}{闭包}。按构造，$K_\eta$ 是\Term{unit-group}{单位群}；又由于 $K$ 在 $S$ 上\Term{separated}{分离}（因为 $H$ 如此），其\Term{unit-section}{单位截面}是闭的，从而 $K$ 是\Term{unit-group}{单位群}。所以 $u$ 是\Term{monomorphism}{单态射}；既然已经看到它\Term{faithfully-flat}{忠实平坦}，它就是\Term{isomorphism}{同构}（参见 SGA 1，I 5.1 或 EGA IV\textsubscript{4}，17.9.1）。这证明 $H$ 是\Term{torus}{环面}，从而完成了 8.2 的证明。\footnote{编者注（65）：至少在至此所考虑的情形，即 $S$ \Term{locally-noetherian}{局部诺特}时如此。}

\textbf{注记 8.5.1。} 也可不引用 8.5，而引用 Zariski“\Term{zariski-main-theorem}{主定理}”；它直接蕴涵 $u$ 是\Term{open-immersion}{开浸入}，从而是\Term{isomorphism}{同构}。\footnote{编者注（66）：编者未能理解此注记，因为不明白为何先验地 $u$ 的\Term{fiber}{纤维}应是\Term{finite}{有限}的，且 $u$ 应为\Term{surjective}{满射}……}

要证明 8.1，由\Term{quasi-isotrivial}{拟等平凡}性\Term{theorem}{定理}可归结为 $S$ 为\Term{local-scheme}{局部}的、$s$ 为其\Term{closed-point}{闭点}的情形，\footnote{编者注（67）：详细说明这一步归约……} 并证明在其余假设下，$H$ 是\Term{torus}{环面}。由已证的 8.2，只需证明 $H$ 的\Term{fiber}{纤维}均为\Term{torus}{环面}。可假定 $S$ 是\Term{complete}{完备}\Term{local-noetherian-ring}{局部诺特环} $A$ 的\Term{spectrum}{谱}。由 3.3，对每个 $n>0$，存在一个在 $S$ 上\Term{finite}{有限}的\Term{multiplicative-type}{乘法型群} $T_n$，以及\Term{isomorphism}{同构} $(T_n)_s\simeq{}_n(H_s)$，其中 $s$ 是 $S$ 的\Term{closed-point}{闭点}。
\SourcePage{97}
像 3.1 中那样进行，并利用 $T_n$ 在 $S$ 上\Term{finite}{有限}，\footnote{编者注（68）：详细说明这一点……} 可见上述\Term{isomorphism}{同构}来自一个\Term{homomorphism}{同态} $u_n:T_n\to H$，且此\Term{homomorphism}{同态}唯一确定。（对 $A$ 的\Term{artinian-ring}{阿廷}\Term{quotient-ring}{商环}取\Term{limit}{极限}。）
此外，由唯一性，当 $m$ 为 $n$ 的倍数时，$u_n:T_n\to H$ 由 $u_m:T_m\to H$ \Term{restriction}{限制}在 ${}_n(T_m)\simeq T_n$ 上得到。由 IX，6.6，$u_n$ 都是\Term{monomorphism}{单态射}，所以各 $T_n$ 是 $H$ 的\Term{subgroup}{子群}；当 $m$ 为 $n$ 的倍数时，有 $T_n={}_n(T_m)$。
所以，对任意 $t\in S$，$(T_n)_t$ 是 $H_t$ 中型为 $(\mathbb Z/n\mathbb Z)^r$ 的\Term{subgroup}{子群}（其中 $r=\dim H_s=\dim H_t$），且当 $m$ 为 $n$ 的倍数时，有 $(T_n)_t={}_n(T_m)_t$。$H_t$ 是\Term{torus}{环面}这一事实，现在由下面的\Term{lemma}{引理}得到。
\begin{SGAresult}{\Term{lemma}{引理}}{8.6}
设 $H$ 是\Term{algebraically-closed}{代数闭}域 $k$ 上的\Term{affine}{仿射}\Term{smooth}{光滑}\Term{algebraic-group}{代数群}，$(T_n)_{n>0}$ 是 $H$ 的一族\Term{multiplicative-type}{乘法型}\Term{subgroup}{子群}，满足：对每个整数 $n>0$，$T_n$ 为型 $(\mathbb Z/n\mathbb Z)^r$，且对 $n$ 的每个倍数 $m$，有 $T_n={}_n(T_m)$。
在这些条件下，$H$ 包含一个\Term{dimension}{维数} $\geq r$ 且包含各 $T_n$ 的\Term{torus}{环面}；所以若 $H$ \Term{connected}{连通}且\Term{dimension}{维数} $\leq r$，则 $H$ 是\Term{dimension}{维数}为 $r$ 的\Term{torus}{环面}。
\end{SGAresult}
\begin{proof}
这是一个\Term{affine}{仿射}\Term{algebraic-group}{代数群}的习题，我们将借助对《Bible》的若干引用来处理。只考虑 $n$ 与\Term{characteristic}{特征}互素的 $T_n$。设 $K$ 是各 $T_n$ 的并在 $H$ 中的\Term{closure}{闭包}，配备诱导的\Term{reduced}{约化}结构；标准论证表明，$K$ 是 $H$ 的一个\Term{commutative-group}{交换}\Term{algebraic-subgroup}{代数子群}。由《Bible》，第 4.5 节，\Term{theorem}{定理} 4，$K$ \Term{isomorphism}{同构}于乘积 $K_u\times K_s$，其中 $K_u$“\Term{unipotent}{幂幺}”，$K_s$ \Term{diagonalizable}{可对角化}。$K$ 的每个\Term{diagonalizable}{可对角化}\Term{subgroup}{子群}都包含在 $K_s$ 中，因而各 $T_n$ 是 $K_s$ 的\Term{subgroup}{子群}，所以 $K=K_s$。写 $K=D(M)$，其中 $M$ 是在 $\mathbb Z$ 上\Term{finitely-generated}{有限生成}的普通\Term{commutative-group}{交换群}；则 $T_n\subseteq K$ 意味着 $M$ 有一个\Term{isomorphism}{同构}于 $(\mathbb Z/n\mathbb Z)^r$ 的\Term{quotient-group}{商群}。这对每个与 $k$ 的\Term{characteristic}{特征}互素的整数 $n$ 都成立（实际上只需对某个固定素数的各次幂成立），因而 $M$ 的\Term{rank}{秩} $\geq r$，所以 $K$ 包含一个\Term{dimension}{维数}为 $r$ 的\Term{torus}{环面}，记作 $T$。当 $H$ \Term{connected}{连通}且\Term{dimension}{维数}为 $r$ 时，得到 $T=H$，这完成了 8.6 的证明。因此 8.1 得证。
\end{proof}

\textbf{注记 8.7。} 利用 8.1，应该不难给出 4.7 和 4.8 的类似推广。更值得研究的是在 8.1 的情形中去掉 $H$ 的\Term{fiber}{纤维}\Term{affine}{仿射}的假设。可以证明，此时存在 $s$ 的一个\Term{open-neighborhood}{开邻域} $U$，使 $H|U$ \Term{commutative-group}{交换}，并且对每个 $t\in U$，\Term{geometric-fiber}{几何纤维} $H_{\overline t}$ 是一个\Term{abelian-variety}{阿贝尔簇}被一个\Term{torus}{环面}的\Term{extension}{扩张}。\footnote{编者注（69）：此处给一个参考文献？} 当然，在这类问题中，只需考虑 $S$ 为\Term{discrete-valuation-ring}{离散赋值环}的\Term{spectrum}{谱}、$s$ 为其\Term{closed-point}{闭点}而 $t$ 为其\Term{generic-point}{泛点}的情形。

可以如下推广这一结果。对于\Term{algebraically-closed}{代数闭}域 $k$ 上的任意\Term{connected}{连通}\Term{smooth}{光滑}\Term{algebraic-group}{代数群} $G$，Chevalley 的一个著名\Term{theorem}{定理}
\SourcePage{98}
告诉我们，$G$（唯一地）是一个\Term{abelian-variety}{阿贝尔簇} $A$ 被一个\Term{affine}{仿射}\Term{smooth}{光滑}\Term{connected}{连通}群 $V$ 的\Term{extension}{扩张}。将 $A$ 的\Term{dimension}{维数}、$G$ 的极大\Term{torus}{环面}的\Term{dimension}{维数}、$G$ 的 \Term{cartan-subgroup}{Cartan 子群}\footnote{编者注（70）：回顾“\Term{cartan-subgroup}{Cartan 子群}”的定义……} 的\Term{dimension}{维数}、$G$（或 $V$）对其\Term{radical}{根基}的商的\Term{dimension}{维数}，分别称为 $G$ 的\Term{abelian-rank}{阿贝尔秩}、\Term{reductive-rank}{约化秩}、\Term{nilpotent-rank}{幂零秩}、\Term{semisimple-dimension}{半单维数}，分别记为 $\rho_{\mathrm{ab}}(G)$、$\rho_r(G)$、$\rho_n(G)$、$d_s(G)$（所有这些概念见《Bible》）。再引入\Term{unipotent-rank}{幂幺秩} $\rho_u(G)=\rho_u(V)=\rho_n(G)-\rho_r(G)-\rho_{\mathrm{ab}}(G)$。当 $G$ 并非\Term{algebraically-closed}{代数闭}时，我们仍用同样的名称和记号表示 $G_{\overline k}$ 的相应\Term{invariants}{不变量}，其中 $\overline k$ 是 $k$ 的\Term{algebraic-closure}{代数闭包}。

现在设 $G$ 是\Term{discrete-valuation-ring}{离散赋值环}的\Term{spectrum}{谱} $S$ 上的\Term{smooth}{光滑}\Term{group-scheme}{群概形}，$\rho_{\mathrm{ab}}$ 等（分别地，$\rho'_{\mathrm{ab}}$ 等）为\Term{special-fiber}{特殊纤维}（分别地，\Term{generic-fiber}{泛纤维}）的上述\Term{invariants}{不变量}。则有不等式
\[
\begin{cases}
\rho_{\mathrm{ab}}\leq\rho'_{\mathrm{ab}},\\
\rho_r+\rho_{\mathrm{ab}}\leq\rho'_r+\rho'_{\mathrm{ab}},\\
d_s\leq d'_s,
\end{cases}
\qquad
\begin{cases}
\rho_n\geq\rho'_n,\\
\rho_u\geq\rho'_u.
\end{cases}
\]
\Term{equivalence}{等价}地说，若 $G$ 在任意底 $S$ 上\Term{smooth}{光滑}\Term{finite-type}{有限型}，则函数 $s\mapsto\rho_{\mathrm{ab}}(s)$、$\rho_{\mathrm{ab}}(s)+\rho_r(s)$、$d_s(s)$ \Term{lower-semicontinuous}{下半连续}，而函数 $s\mapsto\rho_n(s)$、$\rho_u(s)$ \Term{upper-semicontinuous}{上半连续}。\footnote{编者注（71）：为这些结果给一个参考文献？}

即使不假定 $G$ 在 $S$ 上\Term{smooth}{光滑}，只假定它在 $S$ 上\Term{flat}{平坦}\Term{finite-presentation}{有限表示}，这些结果大概仍然成立；此时约定，对\Term{algebraically-closed}{代数闭}域 $k$ 上的\Term{algebraic-group}{代数群} $G$，用 $\rho_{\mathrm{ab}}(G)$ 等表示 $G^0_{\mathrm{r\text{é}d}}$ 的相应\Term{invariants}{不变量}。

在本研讨班中，我们对 $G$ 在 $S$ 上\Term{affine}{仿射}、或更一般地\Term{fiber}{纤维}\Term{affine}{仿射}的情形给出这类结果：在此情形，将验证 $\rho_r$、$\rho_n$，因而 $\rho_u=\rho_n-\rho_r$ 的\Term{semicontinuity}{半连续性}，以及 $\rho_s$ 在\Term{fiber}{纤维}为\Term{reductive-group}{约化群}的点 $s$ 附近的\Term{continuity}{连续性}。\footnote{编者注（71）：为这些结果给一个参考文献？}

若已假定 $G$ \Term{commutative-group}{交换}，可如下推广 8.2：
\begin{SGAresult}{\Term{theorem}{定理}}{8.8}
\footnote{编者注（72）：G. Prasad 与 J.-K. Yu（在若干附加假设下）推广了这一结果，不再假定 $G$ \Term{commutative-group}{交换}，并将假设 a) 及结论中的“\Term{torus}{环面}”换成“\Term{reductive-group}{约化群}”，参见 [PY06]，\Term{theorem}{定理} 6.2。} 设 $G$ 是 $S$ 上\Term{flat}{平坦}\Term{finite-presentation}{有限表示}的\Term{commutative-group}{交换}\Term{group-prescheme}{群预概形}，其\Term{fiber}{纤维}\Term{affine}{仿射}且\Term{connected}{连通}。设 $s\in S$，并假定：

a) 若 $\overline k$ 表示 $\kappa(s)$ 的一个\Term{algebraic-closure}{代数闭包}，则 $(G_{\overline k})_{\mathrm{r\text{é}d}}$ 是\Term{torus}{环面}。

b) 存在 $s$ 的一个\Term{generalization}{一般化} $t$，使 $G_t$ 在 $\kappa(t)$ 上\Term{smooth}{光滑}。

在这些条件下，存在 $s$ 的一个\Term{open-neighborhood}{开邻域} $U$，使 $G|U$ 是\Term{torus}{环面}。
\end{SGAresult}
（另须注意，若只假定对 $s$ 的每个\Term{generalization}{一般化} $t$，$G_t$ \Term{affine}{仿射}，分别地\Term{connected}{连通}，也容易由此\Term{pushout}{推出}：当 $t$ 遍历 $s$ 的一个\Term{open-neighborhood}{开邻域}时，同样的结论成立。）
\SourcePage{99}
\begin{proof}[8.8 的证明]
只需证明 $G_s$ 在 $\kappa(s)$ 上\Term{smooth}{光滑}。事实上，因 $G$ 在 $S$ 上\Term{flat}{平坦}\Term{finite-presentation}{有限表示}，这就\Term{pushout}{推出} $G$ 在 $s$ 的某个\Term{neighborhood}{邻域}上方相对于 $S$ \Term{smooth}{光滑}（参见 3.5），而此时便满足 8.1 的条件。

为证明 $G_s$ 在 $\kappa(s)$ 上\Term{smooth}{光滑}，用通常的方法可归结为 $S$ \Term{affine}{仿射}\Term{noetherian}{诺特}的情形。选取从\Term{discrete-valuation-ring}{离散赋值环}的\Term{spectrum}{谱} $S'$ 到 $S$ 的一个\Term{homomorphism}{同态} $S'\to S$，使它将\Term{closed-point}{闭点}映到 $s$，将\Term{generic-point}{泛点} $\eta$ 映到 $t$；便归结为 $S$ 本身是一个\Term{discrete-valuation-ring}{离散赋值环}的\Term{spectrum}{谱}的情形；还可假定该环\Term{complete}{完备}且\Term{residue-field}{剩余域}\Term{algebraically-closed}{代数闭}，而 $s$ 与 $t=\eta$ 分别是 $S$ 的\Term{closed-point}{闭点}与\Term{generic-point}{泛点}。

于是 $G$ 在 $S$ 上\Term{flat}{平坦}、\Term{separated}{分离}、\Term{finite-type}{有限型}，\Term{generic-fiber}{泛纤维} $G_\eta$ \Term{smooth}{光滑}且\Term{connected}{连通}，而\Term{special-fiber}{特殊纤维} $G_0$ 满足 $T_0=(G_0)_{\mathrm{r\text{é}d}}$ 是\Term{torus}{环面}。设整数 $m>0$ 与 $s$ 处的\Term{residue-characteristic}{剩余特征}互素，因而也与 $\eta$ 处的\Term{residue-characteristic}{剩余特征}互素。此时知道 $m\cdot\operatorname{id}_G$ 逐\Term{fiber}{纤维}为\Term{etale}{平展}\Term{morphism}{态射}（VII\textsubscript{A} 8.4），因 $G$ 在 $S$ 上\Term{flat}{平坦}，它便是\Term{etale}{平展}\Term{morphism}{态射}（SGA 1 I 5.9）；所以其\Term{kernel}{核} ${}_mG$ 在 $S$ 上\Term{etale}{平展}；又因 $G$ 在 $S$ 上\Term{separated}{分离}\footnote{原注（$*$）：由 Raynaud \Term{theorem}{定理} VI\textsubscript{B} 5.3。}\footnote{编者注（73）：参见 8.5 中的编者注（62）。} 且\Term{finite-type}{有限型}，${}_mG$ 也如此。由于 $({}_mG)_0={}_m(T_0)$，其\Term{degree}{次数}为 $m^r$，其中 $r=\dim T_0=\dim G_0=\dim G$。因而 $({}_mG)_\eta={}_m(G_\eta)$ 的\Term{rank}{秩} $\geq m^r$；这已经证明了（使用 8.6）$G_\eta$ 是\Term{dimension}{维数}为 $r$ 的\Term{torus}{环面}，因为 ${}_mG$ 的两个\Term{fiber}{纤维}有相同的\Term{rank}{秩}，所以如 8.5 中一样，${}_mG$ 在 $S$ 上\Term{finite}{有限}。\footnote{编者注（74）：重新检查前一句……}

注意，因为 $S$ \Term{complete}{完备}且\Term{residue-field}{剩余域}\Term{algebraically-closed}{代数闭}，\Term{finite-etale}{有限平展}\Term{cover}{覆盖} ${}_mG$ \Term{completely-decomposed}{完全分解}，所以 ${}_mG_0$ 的每个点都有 $G$ 在 $S$ 上的一个\Term{section}{截面}经过；特别地，$G$ 在 $S$ 上的某个\Term{section}{截面}经过的 $G_0$ 中的点所成\Term{set}{集合}处处\Term{dense}{稠密}。而有如下\Term{lemma}{引理}。
\end{proof}
\begin{SGAresult}{\Term{lemma}{引理}}{8.9}
设 $S$ 是\Term{locally-noetherian}{局部诺特}\Term{regular}{正则}\Term{dimension}{维数} $1$ 的\Term{prescheme}{预概形}，$G$ 是 $S$ 上\Term{flat}{平坦}\Term{locally-finite-type}{局部有限型}的\Term{group-prescheme}{群预概形}，对 $S$ 的每个\Term{maximal-point}{极大点} $\eta$，$G_\eta$ 都在 $\kappa(\eta)$ 上\Term{smooth}{光滑}（这蕴涵 $G$ \Term{reduced}{约化}）。假定 $G$ 的\Term{normalization}{正规化}\Term{scheme}{概形} $X$ 在 $G$ 上\Term{finite}{有限}（由 Nagata 的结果，若 $S$ 是一个\Term{complete}{完备}\Term{discrete-valuation-ring}{离散赋值环}的\Term{spectrum}{谱}，就满足此条件，参见 EGA IV\textsubscript{2}，7.7.4），并设 $G'$ 是由 $X$ 在其处相对于 $S$ \Term{smooth}{光滑}的点构成的\Term{open-subset}{开子集}。采用这些记号，有：

a) 若\Term{projection}{投影} $G'\to S$ 为\Term{surjective}{满射}，则 $G'$ 上存在唯一的 $S$-\Term{group-prescheme}{群预概形}结构，使\Term{canonical}{典范}\Term{morphism}{态射} $G'\to G$ 是群\Term{homomorphism}{同态}。

b) 假定对 $S$ 的每个\Term{closed-point}{闭点} $s$，$G_s^0$ 中有\Term{etale-quasisection}{平展拟截面}经过的点所成\Term{set}{集合}，对于 Zariski \Term{topology}{拓扑}在 $G_s^0$ 中\Term{dense}{稠密}。则 $X$ \Term{regular}{正则}，a) 的条件成立，且映射 $\Gamma(G'/S)\to\Gamma(G/S)$ 是\Term{bijective}{双射}。
\end{SGAresult}
在我们关心的情形，此\Term{lemma}{引理}适用，给出一个群\Term{homomorphism}{同态} $u:G'\to G$，其中 $G'$ 在 $S$ 上\Term{smooth}{光滑}，且 $u_\eta$ 是一个\Term{isomorphism}{同构}
\SourcePage{100}
$G'_\eta\xrightarrow\sim G_\eta$。把 $G'$ 换成一个\Term{open-subgroup}{开子群}，可假定 $G'_0$ \Term{connected}{连通}；而 $u_0:G'_0\to G_0$ 诱导一个\Term{surjective}{满射}\Term{morphism}{态射} $G'_0\to T_0$，其中 $T_0$ 是与 $G'_0$ 同\Term{dimension}{维数}的\Term{torus}{环面}，所以容易\Term{pushout}{推出} $G'_0$ 是\Term{torus}{环面}（例如用 8.6，或引用《Bible》，第 7.3 节，\Term{theorem}{定理} 3 a)）。于是由 8.2，$G'$ 是\Term{torus}{环面}；但这时由 IX 6.8，$\operatorname{Ker}u$ 是 $G'$ 的一个\Term{multiplicative-type}{乘法型}\Term{subgroup}{子群}；其\Term{generic-fiber}{泛纤维}归结为\Term{unit-group}{单位群}，所以它就是\Term{unit-group}{单位群}，因而 $u$ 是\Term{monomorphism}{单态射}。现在使用 VIII 7.9，得到 $u$ 是\Term{immersion}{浸入}。由于它为\Term{surjective}{满射}且 $G$ \Term{reduced}{约化}，$u$ 就是\Term{isomorphism}{同构}，这完成了 8.8 的证明。

\begin{proof}[8.9 的证明]
只剩证明 8.9。为证明 a)，先注意使 $u:G'\to G$ 成为群\Term{homomorphism}{同态}的 $G'$ 上群律的唯一性是显然的，因为已知 $G'$ 在\Term{generic-fiber}{泛纤维}上的群律（假定 $S$ \Term{irreducible}{不可约}，这是允许的）。关于存在性，容易归结为 $S$ 是\Term{local-scheme}{局部}的、为\Term{discrete-valuation-ring}{离散赋值环} $A$ 的\Term{spectrum}{谱}的情形；由唯一性，以及\Term{integral-closure}{整闭包}运算与底的\Term{etale}{平展}\Term{group-extension}{扩张}相交换这一事实，可对 $S$ 作\Term{etale}{平展}\Term{base-change}{基变换}，从而归结为 $A$“\Term{strictly-local}{严格局部}”的情形，即 $A$ 为 \Term{henselian-ring}{Hensel 环}且\Term{residue-field}{剩余域}\Term{separably-closed}{可分闭}。同一归约也适用于 b)，但在 b) 的假设下，现在可将“\Term{etale-quasisection}{平展拟截面}”换成“\Term{section}{截面}”。

注意，由于 $G'$ 在 $S$ 上\Term{smooth}{光滑}且 $X$ \Term{normal}{正规}，$G'\times_SX$ \Term{normal}{正规}（因为它在\Term{normal}{正规}的 $X$ 上\Term{smooth}{光滑}）；所以复合\Term{morphism}{态射} $G'\times_SX\to G\times_SG\to G$ \Term{resolution}{分解}经过
\[p:G'\times_SX\longrightarrow X.\]
证明此\Term{morphism}{态射}在 $G'\times_SX$ 的\Term{open-subset}{开子集} $G'\times_SG'$ 上诱导一个\Term{morphism}{态射}
\[G'\times_SG'\longrightarrow G'.\]
因此需证明 $G'_0\times_kG'_0$ 被映入 $X_0$ 的\Term{open-subset}{开子集} $G'_0$；只需验证对每个取值于 $k$ 的 $G'_0$ 的点 $g'_0$，从 $G'_0$ 到 $X_0$ 的\Term{morphism}{态射}
\[h'_0\longmapsto p_0(g'_0,h'_0)\tag{x}\]
的值落在 $G'_0$ 中。因 $G'$ 在 $S$ 上\Term{smooth}{光滑}且 $S$ 为 Hensel 的，每个上述 $g'_0$ 都由 $G'$ 在 $S$ 上的一个\Term{section}{截面} $g'$ 诱导。立即看出，上述\Term{morphism}{态射} (x) 由从 $G'$ 到 $X$ 的\Term{morphism}{态射} $h'\mapsto p(g',h')$ 诱导；后者则由 $G$ 的\Term{automorphism}{自同构} $h\mapsto g\cdot h$ 通过结构转移而来，此\Term{automorphism}{自同构}是由 $g'$ 在 $G$ 中的像\Term{section}{截面} $g$ 给出的\Term{left-translation}{左平移}。所以 $h'\mapsto p(g',h')$ 本身是 $X$ 的\Term{automorphism}{自同构}，因而将 $G'$ 映入 $G'$，这就证明了我们的断言。

还需证明这样得到的 $G'$ 上的\Term{composition-law}{合成律}是群律。结合律立即由\Term{generic-fiber}{泛纤维}（\Term{isomorphism}{同构}于 $G$ 的\Term{generic-fiber}{泛纤维}）上的结合律得出。另一方面，$S$-\Term{prescheme}{预概形} $G$ 的对称\Term{automorphism}{自同构}诱导 $X$ 的一个\Term{automorphism}{自同构}，它保持 $G'$ \Term{stable}{稳定}，因而诱导 $G'$ 的一个\Term{automorphism}{自同构} $\sigma$。可见它对 $G'$ 的\Term{composition-law}{合成律}具有求逆的性质，因为这仍归结为验证某些涉及 $G'$ 在 $S$ 上的\Term{fiber-product}{纤维幂}的图
\SourcePage{101}
是\Term{commutative-diagram}{交换图}；而这些\Term{fiber-product}{纤维幂}在 $S$ 上\Term{smooth}{光滑}，所以只需在\Term{generic-fiber}{泛纤维}上验证其交换性，这是显然的。这证明了 8.9 的 a) 部分。

证明 b)。设 $Z'$ 是满足 $\mathcal O_{X,x}$ 非\Term{regular-ring}{正则}的 $x\in X$ 所成\Term{set}{集合}；由 Nagata 的一个\Term{theorem}{定理}（EGA IV\textsubscript{2}，6.12.6），它是\Term{closed-subset}{闭子集}，显然包含在 $X_0$ 中。设 $Z$ 是它在 $G$ 中的像，因而是 $G_0$ 的一个\Term{closed-subset}{闭子集}。$Z$ 是 $G_0$ 的一个\Term{rare-subset}{稀疏子集}，即不包含 $G_0$ 的任何\Term{maximal-point}{极大点} $y$。事实上，因为 $G_0$ 在 $G$ 中由一个方程 $t=0$ 定义（其中 $t$ 是\Term{discrete-valuation-ring}{离散赋值环} $A$ 的\Term{uniformizer}{一致化参数}），由\Term{principal-ideal-theorem}{主理想定理}，$\mathcal O_{G,y}$ 的\Term{dimension}{维数}为 $1$，所以对 $X$ 中位于 $y$ 上方的每个 $x$，$\mathcal O_{X,x}$ 的\Term{dimension}{维数}为 $1$；因 $X$ \Term{normal}{正规}，从而在\Term{codimension}{余维} $1$ 处\Term{regular}{正则}，$\mathcal O_{X,x}$ 便是\Term{discrete-valuation-ring}{离散赋值环}。另一方面，显然对 $G$ 在 $S$ 上的每个\Term{section}{截面} $g$，$Z'$ 在由 $G$ 中按 $g$ \Term{left-translation}{左平移}、通过结构转移所定义的 $X$ 的\Term{automorphism}{自同构}下\Term{stable}{稳定}，所以 $Z$ 在 $G_0$ 中按 $g_0$ \Term{left-translation}{左平移}下\Term{stable}{稳定}。按假设，这些 $g_0$ 所成\Term{set}{集合}在 $G_0^0$ 中\Term{dense}{稠密}。由于 $Z$ 在这些 $g_0$ 的\Term{translation}{平移}下\Term{stable}{稳定}，且为闭的\Term{rare-subset}{稀疏子集}，立即得到 $Z=\varnothing$，从而 $Z'=\varnothing$，所以 $X$ \Term{regular}{正则}。因此，$X$ 在每个有\Term{section}{截面}经过的点处相对于 $S$ \Term{smooth}{光滑}。而 $G$ 在 $S$ 上的每个\Term{section}{截面}都唯一\Term{lifting}{提升}为 $X$ 在 $S$ 上的一个\Term{section}{截面}，从而为 $G'$ 在 $S$ 上的一个\Term{section}{截面}。特别地，这适用于\Term{unit-section}{单位截面}，证明 $G'$ 在 $S$ 中的像是 $S$，即 a) 的条件成立。这完成了 8.9，从而也完成了 8.8 的证明。
\end{proof}

\SGAsection{9}{增补}
\footnote{编者注（75）：本增补节由 Fabrice Orgogozo 撰写（遵循 Ofer Gabber 的提示）。}
\textbf{9.1。单纯集、\Term{topos}{拓扑斯}、\Term{groupoid}{群胚}与\Term{topological-space}{拓扑空间}。}
\begin{SGAresult}{记号}{9.1.1}
(1) 设 $E$ 是一个\Term{simplicial-set}{单纯集}。可将下列\Term{object}{对象}与它相联系：

(2) 一个\Term{topos}{拓扑斯} $T=E^\sim$，它由自然地与各\Term{set}{集合} $E_i$ 相联系的\Term{topos}{拓扑斯}，按照 [Del74]，6.3.1 所述方法得到（另见 [Ill72]，VI.5.2 和 SGA 4 VI.7）；

(3) 一个\Term{groupoid}{群胚} $G=\Pi(E)$，其\Term{object}{对象}是 $E_0$ 的元素（“顶点”），其\Term{arrow}{箭头}在 [GZ67]，II.7 中定义；

(4) 一个\Term{topological-space}{拓扑空间} $X=|E|$（一个\Term{cell-complex}{胞腔复形}），称为“\Term{geometric-realization}{几何实现}”（或“\Term{topology}{拓扑}实现”）（参见该处，III.1）。
\end{SGAresult}
注意，$T$ 上的一个\Term{sheaf}{层}无非就是 $E$ 上方的一个\Term{simplicial-set}{单纯集}。
\SourcePage{102}
\textbf{9.2。\Term{locally-constant-sheaf}{局部常值层}；\Term{descent-datum}{下降数据}。}
\begin{SGAresult}{定义}{9.2.1}
下列\Term{object}{对象}称为相应\Term{object}{对象}上的\Term{locally-constant-sheaf}{局部常值层}：

(1) 对\Term{simplicial-set}{单纯集} $E$，是任意一个\Term{simplicial-set}{单纯集}\Term{morphism}{态射} $E'\to E$，使对每个 $i\in\mathbb N$ 及每个 $e\in E_i$，\Term{face-operator}{面算子} $d$ 都在\Term{fiber}{纤维}间诱导\Term{isomorphism}{同构} $E'_e\xrightarrow\sim E'_{d(e)}$（参见 [AM69]，第 10 节）；

(2) 对\Term{topos}{拓扑斯} $T$，是 $T$ 的任意\Term{object}{对象} $F$，使存在一个\Term{epimorphism}{满态射} $U\to1$ 和一个\Term{isomorphism}{同构} $F\times U\simeq f^*I\times U$，其中 $I$ 为一个\Term{set}{集合}，$f:T\to\mathrm{Ens}$ 为终\Term{morphism}{态射}（参见 SGA 4，IX.2）；

(3) 对\Term{groupoid}{群胚} $G$，是 $G$ 上的任意\Term{presheaf}{预层}，即从 $G$ 到\Term{set}{集合}\Term{category}{范畴}的任意\Term{contravariant-functor}{反变函子}（参见 [GZ67]，附录 I.1.2）；

(4) 对\Term{topological-space}{拓扑空间} $X$，是 $X$ 上在通常意义下\Term{locally-constant}{局部常值}的任意\Term{set}{集合}\Term{sheaf}{层}。

最后：

(5) 称\Term{simplicial-set}{单纯集} $E$ 上的一个\Term{descent-datum}{下降数据}为如下数据：$E_0^\sim$ 上的一个\Term{sheaf}{层} $F$（即 $E_0$ 上的一个\Term{set}{集合}值函数），以及一个\Term{isomorphism}{同构} $\alpha:p_1^*F\xrightarrow\sim p_2^*F$，其中 $p_1,p_2:E_1^\sim\to E_0^\sim$ 是面（所\Term{derive}{导出}的）\Term{morphism}{态射}，且此\Term{isomorphism}{同构}满足通常的\Term{cocycle}{余圈}关系（参见第 IV 讲，2.1\footnote{PP：我把“FGA，\Term{descent}{下降}技术 I，1.6”替换为“第 IV 讲，2.1”。}，以及下文）。
\end{SGAresult}
这五类\Term{object}{对象}之间的\Term{morphism}{态射}以及相应的\Term{inverse-image-functor}{逆像函子}，均以显然的方式定义。

\textbf{9.3。若干\Term{category-equivalence}{范畴等价}。} 设 $E$ 为\Term{simplicial-set}{单纯集}，$T$、$G$、$X$ 分别为与它相联系的\Term{topos}{拓扑斯}、\Term{groupoid}{群胚}和\Term{topological-space}{拓扑空间}。
\begin{SGAresult}{命题}{9.3.1}
$E$、$T$、$G$、$X$ 上的\Term{locally-constant-sheaf}{局部常值层}的\Term{category}{范畴}，以及 $E$ 上\Term{descent-datum}{下降数据}的\Term{category}{范畴}，彼此\Term{equivalence}{等价}。
\end{SGAresult}
\begin{proof}[证明概要]
用 (1) 至 (5) 表示上一段所定义\Term{object}{对象}的\Term{category}{范畴}。

— (1)$\Leftrightarrow$(5)。这是 [AM69]，10.6 的一个特例（另见 [GZ67]，附录 I.2.3，[Fri82]，\SGAsourcegap{原 PDF 第 102 页此引注印作“?.5.6”，缺少 [Fri82] 的确切章号}，或 [Ill72]，VI.8.1.6）。把\Term{object}{对象} $(E'\to E)$ 映到二元组 $(E'_0,\alpha)$ 的\Term{functor}{函子}给出一个\Term{category}{范畴}\Term{equivalence}{等价}；其中 $E'_0$ 视为 $E_0$ 上的\Term{sheaf}{层}，而 $\alpha$ 是唯一的\Term{isomorphism}{同构} $p_1^*E'_0\xrightarrow\sim p_2^*E'_0$，其在每个 $y\in E_1$ 处的\Term{fiber}{纤维}——将 $y$ 在两个\Term{projection}{投影}下的像记为 $x_1$ 与 $x_2$——由\Term{isomorphism}{同构}
\[(E'_0)_{x_1}\xleftarrow\sim(E'_1)_y\xrightarrow\sim(E'_0)_{x_2}\]
给出。

— (5)$\Leftrightarrow$(3)。显然：定义与 $E$ 相联系的\Term{groupoid}{群胚} $G$ 中\Term{morphism}{态射}的两个关系之一就是一个\Term{cocycle}{余圈}关系。

— (1)$\Leftrightarrow$(4)。参见 [GZ67]，附录 I.3.2.1。

— (1)$\Leftrightarrow$(2)。须证明，$E$ 上方的一个\Term{object}{对象}在\Term{simplicial}{单纯}意义下为\Term{locally-constant-sheaf}{局部常值层}，当且仅当它作为 $T=E^\sim$ 上的\Term{sheaf}{层}是\Term{locally-constant}{局部常值}的。这由如下事实得到：$T$ 上的\Term{locally-constant-sheaf}{局部常值层} $F_\bullet$ 恰为\Term{cartesian-sheaf}{笛卡尔层}，即\Term{arrow}{箭头} $([n]\to[m])^*F_m\to F_n$ 为\Term{isomorphism}{同构}的\Term{sheaf}{层}。这一点是关于\Term{simplicial-topos}{单纯拓扑斯}的一个一般事实，结合 $E_0^\sim$ 上每个\Term{sheaf}{层}均为\Term{locally-constant}{局部常值}这一事实的特例。
\SourcePage{103}
（另见 [Ill72]，VI.8.1.6。）
\end{proof}
对任意群 $H$，上面的\Term{category}{范畴}\Term{equivalence}{等价}诱导 $H$-\Term{torsor}{挠子}\Term{category}{范畴}之间的\Term{equivalence}{等价}；这里的 $H$-\Term{torsor}{挠子}是带有 $H$ 的\Term{action}{作用}的\Term{locally-constant-sheaf}{局部常值层}，其\Term{fiber}{纤维}\Term{isomorphism}{同构}于以\Term{translation}{平移}\Term{action}{作用}于自身的 $H$。

\textbf{9.4。\Term{fundamental-group}{基本群}与\Term{fundamental-groupoid}{基本群胚}。} 由前面的\Term{category}{范畴}\Term{equivalence}{等价}可知，对任意群 $H$ 与任意\Term{simplicial-set}{单纯集} $E$，$H$-\Term{torsor}{挠子}的\Term{isomorphism-class}{同构类}\Term{set}{集合}
\[
\mathrm H^1(E,H),\quad \mathrm H^1(E^\sim,H),\quad
\pi_0(\operatorname{Hom}(\Pi(E),BH)),\quad \mathrm H^1(|E|,H)
\]
之间有自然\Term{bijective}{双射}。回顾一下，$BH$ 表示与 $H$ 相联系的单点\Term{groupoid}{群胚}，$\pi_0$ 是把一个\Term{category}{范畴}映到其\Term{object}{对象}的\Term{isomorphism-class}{同构类}\Term{set}{集合}的\Term{functor}{函子}。
同样，若 $e$ 是 $E$ 的一个点，上述\Term{equivalence}{等价}在 $e$ 上已\Term{trivialization}{平凡化}的 $H$-\Term{torsor}{挠子}的\Term{isomorphism-class}{同构类}\Term{set}{集合}间诱导\Term{bijective}{双射}；它们分别记作
\[
\mathrm H^1(E\mathrel{\mathrm{rel}}e,H),\quad
\mathrm H^1(E^\sim\mathrel{\mathrm{rel}}e^\sim,H),\quad
\operatorname{Hom}(\pi_1(\Pi(E),e),H),\quad
\mathrm H^1(|E|\mathrel{\mathrm{rel}}|e|,H).
\]
回顾一下，用 $\pi_1(\Pi(E),e)$ 表示群 $\operatorname{Isom}_{\Pi(E)}(e,e)$。当 $H$ 变动时，这些\Term{functor}{函子}在\Term{connected}{连通}情形由一个群\Term{representable}{表示}，记该群为 $\pi_1(E,e)$。群 $\pi_1(E,e)$ \Term{isomorphism}{同构}于 $\pi_1(\Pi(E),e)$ 和 $\pi_1(|E|,|e|)$，因而也\Term{isomorphism}{同构}于 Kan 所定义的\Term{simplicial-set}{单纯集}的\Term{fundamental-group}{基本群}（参见例如 [May67]，16.1 或 [Ill72]，I.2.1.1）。（回顾一下，\Term{set}{集合} $\mathrm H^1(E,H)$ 则\Term{isomorphism}{同构}于到 $H$ 的\Term{morphism}{态射}模去\Term{conjugacy}{共轭}所成\Term{set}{集合} $\mathrm H^1(\pi_1(E,e),H)$，也记作 $\operatorname{Hom\,ext}(\pi_1(E,e),H)$。）

\textbf{9.5。锥。}
\begin{SGAresult}{定义}{9.5.1}
(1) 设 $f:E'\to E$ 是\Term{simplicial-set}{单纯集}\Term{morphism}{态射}。回顾一下（参见例如 [Del74]，6.3.1），用 $C(f)$ 表示如下\Term{pointed-simplicial-set}{带点单纯集}，其\Term{degree}{次数} $n\geq0$ 的\Term{underlying-set}{底集}为
\[
E_n\amalg\Bigl(\coprod_{i<n}E'_i\Bigr)\amalg\star,
\]
其中 $\star$ 是单元素\Term{set}{集合}。我们将各\Term{simplicial}{单纯}映射的定义留给读者，下面会回顾\Term{degree}{次数}不超过二的面映射的定义。（另见 [GZ67]，VI.2，其中给出一个带点的变体。）$C(f)$ 上\Term{locally-constant-sheaf}{局部常值层}的\Term{category}{范畴}，\Term{equivalence}{等价}于 $E$ 上配备其在 $E'$ 上\Term{inverse-image}{逆像}的一个\Term{trivialization}{平凡化}的\Term{locally-constant-sheaf}{局部常值层}的\Term{category}{范畴}。\Term{simplicial-set}{单纯集} $C(f)$ 自然地由 $\star\to C(f)$ 带点。

(2) 设 $f:T'\to T$ 是\Term{topos}{拓扑斯}\Term{morphism}{态射}。用 $C(f)$ 表示如下\Term{topos}{拓扑斯}，其\Term{object}{对象}是五元组
\[
(F,F',A,\alpha:f^*F\to F',\beta:A\to\Gamma(T',F')),
\]
其中 $F$（分别地 $F'$）是 $T$（分别地 $T'$）的\Term{object}{对象}，$A$ 是一个\Term{set}{集合}，而 $\alpha$（分别地 $\beta$）是 $T'$（分别地 $\mathrm{Ens}$）中的\Term{morphism}{态射}。（另见 [Del80]，4.3.4 与 [Ill72]，III.4，其中给出了此构造的一个变体。）$C(f)$ 上\Term{locally-constant-sheaf}{局部常值层}的\Term{category}{范畴}，\Term{equivalence}{等价}于 $T$ 上配备其在 $T'$ 上\Term{inverse-image}{逆像}的一个\Term{trivialization}{平凡化}的\Term{locally-constant-sheaf}{局部常值层}的\Term{category}{范畴}。
\SourcePage{104}
\Term{topos}{拓扑斯} $C(f)$ 自然地由将五元组 $(F,F',A,\alpha:f^*F\to F',\beta:A\to\Gamma(T',F'))$ 映到\Term{set}{集合} $A$ 的\Term{fiber-functor}{纤维函子}赋予基点。

(3) 设 $f:G'\to G$ 是\Term{groupoid}{群胚}\Term{morphism}{态射}。用 $C(f)$ 表示图
\[G\longleftarrow G'\longrightarrow B1\]
的\Term{colimit}{余极限}，其中 $B1$ 是单点\Term{category}{范畴}（一个\Term{object}{对象}，一支\Term{arrow}{箭头}）。$C(f)$ 上\Term{locally-constant-sheaf}{局部常值层}的\Term{category}{范畴}，\Term{equivalence}{等价}于 $G$ 上配备其在 $G'$ 上\Term{inverse-image}{逆像}的一个\Term{trivialization}{平凡化}的\Term{locally-constant-sheaf}{局部常值层}的\Term{category}{范畴}。

(4) 设 $f:X'\to X$ 是\Term{topological-space}{拓扑空间}\Term{morphism}{态射}。用 $C(f)$ 表示图
\[
\star\longleftarrow X'\times\{0\}\longrightarrow X'\times[0,1]
\longleftarrow X'\times\{1\}\longrightarrow X
\]
的\Term{colimit}{余极限}。$C(f)$ 上\Term{locally-constant-sheaf}{局部常值层}的\Term{category}{范畴}，\Term{equivalence}{等价}于 $X$ 上配备其在 $X'$ 上\Term{inverse-image}{逆像}的一个\Term{trivialization}{平凡化}的\Term{locally-constant-sheaf}{局部常值层}的\Term{category}{范畴}。\Term{topological-space}{拓扑空间} $C(f)$ 自然地由 $\star\to C(f)$ 带点。
\end{SGAresult}

\textbf{9.5.2。单纯映射的锥上的\Term{descent-datum}{下降数据}。}
设 $f:E'\to E$ 是\Term{simplicial-set}{单纯集}\Term{morphism}{态射}，$C(f)$ 是其\Term{cone}{锥}（参见上面 (1)）。用字母 $p$（分别地 $q$、$r$）表示 $E'$（分别地 $E$、$C(f)$）的面映射。在陈述下面的命题之前，先用 $p$ 与 $q$ 明确写出\Term{degree}{次数}不超过二的面 $r$。约定 $p_{ij}$（分别地 $p_i$）是对应于像为 $\{i,j\}$（分别地 $\{i\}$）的递增映射 $\{1,2\}\to\{1,2,3\}$（分别地 $\{1\}\to\{1,2\}$）的面映射
\[
E'_2=E'_{\{1,2,3\}}\longrightarrow E'_1=E'_{\{1,2\}}
\quad\text{（分别地 }E'_1=E'_{\{1,2\}}\longrightarrow E'_0=E'_{\{1\}}\text{）。}
\]
\footnote{原注（2）：这一约定与目前的\Term{simplicial}{单纯}记号惯例不同，但更接近 Grothendieck 在其\Term{descent}{下降}理论中使用的记号。} 对 $E$ 与 $C(f)$ 的面采用同样的约定。
\Term{morphism}{态射}
\[
r_1:C(f)_1=E_1\amalg E'_0\amalg\star\longrightarrow C(f)_0=E_0\amalg\star
\]
（分别地 $r_2$）为：

— 在 $E_1$ 上是 $q_1:E_1\to E_0$（分别地 $q_2:E_1\to E_0$）；

— 在 $E'_0$ 上是 $E'_0\to\star$（分别地 $f_0:E'_0\to E_0$）；

— 在 $\star$ 上是 $\star\to\star$。

同样，\Term{morphism}{态射}
\[
r_{21}:C(f)_2=E_2\amalg E'_1\amalg E'_0\amalg\star
\longrightarrow C(f)_1=E_1\amalg E'_0\amalg\star
\]
（分别地 $r_{32}$、$r_{31}$）为：

— 在 $E_2$ 上是 $q_{21}:E_2\to E_1$（分别地 $q_{32}$、$q_{31}$）；

— 在 $E'_1$ 上是 $p_1:E'_1\to E'_0$（分别地 $f_1:E'_1\to E_1$、$p_2:E'_1\to E'_0$）；

— 在 $E'_0$ 上是 $E'_0\to\star$（分别地 $\operatorname{Id}:E'_0\to E'_0$、$\operatorname{Id}:E'_0\to E'_0$）。

设 $H$ 是 $C(f)$ 上的一个\Term{descent-datum}{下降数据}，即 $C(f)_0=E_0\amalg\star$ 上的一个\Term{sheaf}{层} $H_0$，配有其在 $C(f)_1=E_1\amalg E'_0\amalg\star$ 上的\Term{inverse-image}{逆像}之间的一个\Term{isomorphism}{同构} $\gamma:r_1^*H_0\xrightarrow\sim r_2^*H_0$，满足\Term{cocycle}{余圈}关系 $r_{31}^*\gamma=r_{32}^*\gamma\circ r_{21}^*\gamma$。
\SourcePage{105}
\Term{isomorphism}{同构} $\gamma$ 在 $E_1$（分别地 $E'_0$）上的\Term{restriction}{限制}，是一个\Term{descent-datum}{下降数据} $\beta:q_1^*G_0\xrightarrow\sim q_2^*G_0$，其中 $G_0$ 是 $H_0$ 在 $E_0$ 上的\Term{restriction}{限制}（分别地，是一个\Term{trivialization}{平凡化} $t:\underline{H_0(\star)}\xrightarrow\sim f^*G_0=:F_0$，其中 $\underline{H_0(\star)}$ 是以 $H_0(\star)$ 为茎的\Term{constant-sheaf}{常值层}）。$\gamma$ 所满足的\Term{cocycle}{余圈}关系在 $E_2$（分别地 $E'_1$）上的\Term{restriction}{限制}，是 $\beta$ 的\Term{cocycle}{余圈}关系（分别地，是 $p_2^*(t)=f_1^*(\beta)\circ p_1^*(t)$）。后一个关系意味着\Term{trivialization}{平凡化} $t$ 与 $\beta$ 在 $F_0$ 上诱导的\Term{descent-datum}{下降数据} $f_1^*(\beta)$ 相容。因此：
\begin{SGAresult}{命题}{9.5.3}
设 $f:E'\to E$ 是\Term{simplicial-set}{单纯集}\Term{morphism}{态射}，$C(f)$ 是其\Term{cone}{锥}，以 $\star$ 为基点。$C(f)$ 上在 $\star$ 上已\Term{trivialization}{平凡化}的\Term{descent-datum}{下降数据}的\Term{category}{范畴}，\Term{equivalence}{等价}于 $E$ 上配备其在 $E'$ 上诱导的\Term{descent-datum}{下降数据}的一个\Term{trivialization}{平凡化}的\Term{descent-datum}{下降数据}的\Term{category}{范畴}。
\end{SGAresult}

\textbf{9.6。\Term{functor}{函子} $\pi^1(S'/S,\xi;-)$ 的\Term{representability}{可表性}。} 记号同第 90 页，并且从现在起假定 $S_0=S'$ 与 $S_1=S'\times_SS'$ 的\Term{connected-component}{连通分支}都是开的。在此假设下，有如下明确的组合描述，它立即由上面的回顾得到。
\begin{SGAresult}{命题}{9.6.1}
对任意群 $G$，配备 $\xi$ 上方的一个\Term{trivialization}{平凡化}的\Term{descent-datum}{下降数据}的\Term{isomorphism-class}{同构类}\Term{set}{集合} $\pi^1(S'/S,\xi;G)$，与相对\Term{nonabelian-cohomology}{非阿贝尔上同调}\Term{set}{集合}
\[
\mathrm H^1(K\mathrel{\mathrm{rel}}k,G):=\mathrm H^1(\widetilde K\mathrel{\mathrm{rel}}\star,G)
\]
之间有自然\Term{bijective}{双射}。
\end{SGAresult}
由此得到，只要\Term{simplicial-set}{单纯集} $\widetilde K=C(k\to K)$ \Term{connected}{连通}，此\Term{functor}{函子}就由一个群\Term{representable}{表示}，记作 $\pi_1(K,k)$（“\Term{relative-fundamental-group}{相对基本群}”）。立即可以验证，此条件成立当且仅当映射 $\pi_0(k)\to\pi_0(K)$ 为\Term{surjective}{满射}。特别地，只需\Term{simplicial-set}{单纯集} $K$ \Term{connected}{连通}。如正文所指出，若\Term{scheme}{概形} $S$ \Term{connected}{连通}且\Term{morphism}{态射} $S'\to S$ \Term{universally-submersive}{普遍淹没}，就属于此情形。

\textbf{9.7。$k$ 的\Term{contractible}{可缩}性，$k_\bullet\to K_\bullet$ \Term{injectivity}{单射}性的反例。} 下面两个命题详细说明编者注（53）。
\begin{SGAresult}{命题}{9.7.1}
设 $\kappa$ 是\Term{algebraically-closed}{代数闭}域，$X$ 是\Term{connected}{连通}的 $\kappa$-\Term{scheme}{概形}。对每个整数 $i\geq0$，用 $X_i$ 表示 $X$ 在 $\kappa$ 上的第 $(i+1)$ 次\Term{fiber-product}{纤维幂}。对每个整数 $i\geq0$，\Term{canonical}{典范}映射
\[\pi_0(X_i)\longrightarrow\pi_0(X)^{i+1}\]
是\Term{bijective}{双射}。特别地，\Term{simplicial-set}{单纯集} $\pi_0(X_\bullet)$ \Term{contractible}{可缩}。
\end{SGAresult}
只需证明下面的\Term{lemma}{引理}，它由 Künneth 公式通过取\Term{limit}{极限}得到。[详细说明？]
\begin{SGAresult}{\Term{lemma}{引理}}{9.7.2}
设 $\kappa$ 是\Term{algebraically-closed}{代数闭}域，$X$ 与 $Y$ 是两个\Term{connected}{连通}的 $\kappa$-\Term{scheme}{概形}。则\Term{fiber-product}{纤维积} $X\times_\kappa Y$ \Term{connected}{连通}。
\end{SGAresult}
\begin{SGAresult}{命题}{9.7.3}
设 $X$ 是\Term{connected}{连通}\Term{scheme}{概形}，$X'\to X$ 是\Term{finite-etale}{有限平展}\Term{morphism}{态射}，$\overline x$ 是 $X$ 的一个位于点 $x$ 的\Term{geometric-point}{几何点}。\Term{canonical}{典范}\Term{morphism}{态射} $\pi_0(X'_{\overline x})\to\pi_0(X')$ 为\Term{surjective}{满射}。
\end{SGAresult}
\SourcePage{106}
\begin{proof}
事实上，$X'$ 的每个\Term{connected-component}{连通分支}在 $X$ 中的像都是既开又闭的，所以与\Term{fiber}{纤维} $X'_x$ 相交。因此\Term{morphism}{态射} $\pi_0(X'_x)\to\pi_0(X')$ 为\Term{surjective}{满射}，而\Term{morphism}{态射} $\pi_0(X'_{\overline x})\to\pi_0(X'_x)$ 也为\Term{surjective}{满射}。
\end{proof}

\textbf{参考文献}\footnote{编者注（76）：本讲所引的补充参考文献；[AM69] 及后面的参考文献与增补有关。}
\begin{description}
\item[{[BAC]}] N. Bourbaki，\textit{\Term{commutative-algebra}{交换代数}}（\textit{Algèbre commutative}），第 I—IV 章，Hermann，1961；Masson，1985；Springer-Verlag，2006。
\item[{[BEns]}] N. Bourbaki，\textit{\Term{set-theory}{集合论}}（\textit{Théorie des ensembles}），第 I—IV 章，Hermann，1970。
\item[{[Pes66]}] C. Peskine，Zariski“主\Term{theorem}{定理}”的一个推广（\textit{Une généralisation du « Main Theorem » de Zariski}），\textit{Bull. Sci. Math.} \textbf{90}（1966），119—127。
\item[{[PY06]}] G. Prasad、J.-K. Yu，关于拟\Term{reductive-group-scheme}{约化群概形}（\textit{On quasi-reductive group schemes}），\textit{J. Alg. Geom.} \textbf{15}（2006），507—549。
\item[{[Ray70]}] M. Raynaud，\textit{\Term{henselian-local-ring}{Hensel 局部环}}（\textit{Anneaux locaux henséliens}），\textit{Lect. Notes Maths.} \textbf{169}，Springer-Verlag，1970。
\item[{[AM69]}] M. Artin、B. Mazur，\textit{\Term{etale}{平展}\Term{homotopy}{同伦}}（\textit{Etale homotopy}），\textit{Lect. Notes Maths.} \textbf{100}，Springer，1969。
\item[{[Del74]}] P. Deligne，\Term{hodge-theory}{Hodge 理论} III（\textit{Théorie de Hodge III}），\textit{Publ. Math. IHÉS} \textbf{44}（1974），5—77。
\item[{[Del80]}] P. Deligne，Weil 猜想 II（\textit{La conjecture de Weil II}），\textit{Publ. Math. IHÉS} \textbf{52}（1980），137—252。
\item[{[Fri82]}] E. M. Friedlander，\textit{单纯\Term{scheme}{概形}的\Term{etale}{平展}\Term{homotopy}{同伦}}（\textit{Etale homotopy of simplicial schemes}），Princeton Univ. Press，1982。
\item[{[GZ67]}] P. Gabriel、M. Zisman，\textit{\Term{calculus-of-fractions}{分式演算}与\Term{homotopy}{同伦}理论}（\textit{Calculus of fractions and homotopy theory}），Springer，1967。
\item[{[Ill72]}] L. Illusie，\textit{\Term{cotangent-complex}{余切复形}与\Term{deformation}{形变} I、II}（\textit{Complexe cotangent et déformations I \& II}），\textit{Lect. Notes Maths.} \textbf{239}、\textbf{283}，Springer，1971、1972。
\item[{[Kan58]}] D. Kan，\Term{homotopy-group}{同伦群}的一个组合定义（\textit{A combinatorial definition of homotopy groups}），\textit{Ann. of Math.} \textbf{67}（1958），282—312。
\item[{[May67]}] J. P. May，\textit{\Term{algebra}{代数}\Term{topology}{拓扑}中的\Term{simplicial-object}{单纯对象}}（\textit{Simplicial objects in algebraic topology}），Univ. of Chicago Press，1967。
\end{description}
